\documentclass[a4paper,openbib,10pt]{memoir}
\usepackage[utf8]{inputenc}
\usepackage[english]{babel}
\usepackage[T1]{fontenc}
\usepackage{lmodern}
\usepackage{amsmath}
\usepackage{amssymb}
\usepackage{amsthm}
\usepackage{mathtools}
\usepackage{enumerate}
\usepackage{enumitem}
\usepackage[dvips]{graphicx}
\usepackage{wrapfig}
\usepackage{multirow}
\usepackage{fancybox}

\usepackage[dvipsnames]{xcolor}

\colorlet{isomorphism}{gray!75}
\colorlet{injective}{lightgray!75}

\usepackage{listings}
\lstset{
  breakatwhitespace=true,
  breaklines=true,
  frame=single,
  postbreak=\raisebox{0ex}[0ex][0ex]{\ensuremath{\color{red}\hookrightarrow\space}},
  showstringspaces=false,
  stringstyle=\color{green}
}


\usepackage{csquotes}
\usepackage[backend=bibtex,style=alphabetic]{biblatex}
\addbibresource{phdbooks.bib}

\usepackage[english, status=draft]{fixme}

\usepackage[all,cmtip]{xy}
\entrymodifiers={+!!<0pt,\fontdimen22\textfont2>}

\setlrmarginsandblock{*}{4cm}{*}
\checkandfixthelayout

\hyphenation{to-po-lo-gi-cal}

\chapterstyle{southall} 

\theoremstyle{plain}
\newtheorem{theorem}{Theorem}
\newtheorem{lemma}[theorem]{Lemma}
\newtheorem{proposition}[theorem]{Proposition}
\newtheorem{corollary}[theorem]{Corollary}
\newtheorem*{conjecture}{Conjecture}

\theoremstyle{definition}
\newtheorem{definition}[theorem]{Definition}
\newtheorem{example}[theorem]{Example}
\newtheorem*{hypothesis}{Hypothesis}

\theoremstyle{remark}
\newtheorem{remark}[theorem]{Remark}

\newcommand{\ip}[2]{\left< #1, #2 \right>}
\newcommand{\size}[1]{\left| #1 \right|}

\let\O\undefined

\newcommand{\N}{\mathbb{N}}
\newcommand{\Z}{\mathbb{Z}}
\newcommand{\Q}{\mathbb{Q}}
\newcommand{\R}{\mathbb{R}}
\newcommand{\C}{\mathbb{C}}
\newcommand{\A}{\mathcal{A}}
\newcommand{\delim}{\:}
\newcommand{\X}{X_{m,n}}
\newcommand{\Y}{Y_{m,n}}
\newcommand{\sdel}{\delim\Bigg\vert\delim}
\newcommand{\SUT}[1]{\SU_{#1}/\T^{#1-1}}
\newcommand{\GLB}[1]{\GL_{#1}(\C)/\B_{#1}}
\newcommand{\UT}[1]{\U_{#1}/\T^{#1}}
\DeclareMathOperator{\GL}{GL}
\DeclareMathOperator{\U}{U}
\DeclareMathOperator{\SU}{SU}
\DeclareMathOperator{\SL}{SL}
\newcommand{\TT}[2]{\T^{#1}\!\!/\T^{#2}}
\newcommand{\Hq}{\mathbb{H}}
\DeclareMathOperator{\O}{O}
\DeclareMathOperator{\SO}{SO}
\DeclareMathOperator{\GS}{GS}
\DeclareMathOperator{\Ext}{Ext}
\DeclareMathOperator{\Hom}{Hom}
\DeclareMathOperator{\im}{im}
\DeclareMathOperator{\Tr}{Tr}
\DeclareMathOperator{\Id}{Id}
\DeclareMathOperator{\supp}{supp}
\DeclareMathOperator{\pr}{pr}
\DeclareMathOperator{\spa}{span}
\DeclareMathOperator{\inv}{inv}
\DeclareMathOperator{\Det}{Det}
\DeclareMathOperator{\T}{T}
\DeclareMathOperator{\B}{B}

\newcommand{\hmat}{\rule[.5ex]{2.5ex}{0.5pt}}
\newcommand{\matrow}[1]{\hmat & \hspace{-0.75em} #1 & \hspace{-0.75em}
  \hmat}

\newcommand{\closure}[2][3]{{}\mkern#1mu\overline{\mkern-#1mu#2}}

\newcommand{\set}[1]{\left\{ #1 \right\}}
\newcommand{\com}[1]{\left[#1,#1\right]}

\newcommand{\lie}[1]{\left[ #1 \right]}

\renewcommand{\a}{\alpha}
\renewcommand{\b}{\beta}
\newcommand{\g}{\gamma}
\newcommand{\e}{\varepsilon}
\renewcommand{\d}{\delta}
\renewcommand{\l}{\lambda}

\begin{document}

\frontmatter

\begin{titlingpage}
  % forside centreret i forhold til papiret ikke i forhold til
  % marginerne
  \calccentering{\unitlength}
  \begin{adjustwidth*}{\unitlength}{-\unitlength}
    \begin{adjustwidth}{-1cm}{-1cm} 
      \vspace*{1cm}
      \centering
      \scshape
      
      {\Huge
        \begin{center}
          Approximations to the loop
          space of flag manifolds
        \end{center}
        }
      
      \vspace*{4.5\onelineskip}
      
      \includegraphics[width=8cm]{billeder/segla1s.pdf}    
      
      \vspace*{4.5\onelineskip}
      
      {\Large
        \linespread{1.4}\selectfont
        {\huge Simon Stoltze}\\
        \vspace{1em}
        Ph.D. Dissertation\\
        October, 2016 % \today
        \\ 
        \vspace{1em}
        Advisor: Marcel B\"okstedt}
      
      \strut\vfill
      Department of Mathematics\\
      Aarhus University
    \end{adjustwidth}
  \end{adjustwidth*}
  
\end{titlingpage}


%%% Local Variables: 
%%% mode: latex
%%% TeX-master: "main"
%%% End: 

\tableofcontents


% Ny struktur. Baseret på gamle studerendes afhandlinger.

\chapter{Preface}

This dissertation is the outcome of my Ph.D. studies at the Department
of Mathematics, Aarhus University, from February 2013 to October
2016. It contains and expands on the subjects covered in my combined
progress report and Master's thesis, which was finished and defended
in January 2015.

I would like to thank my Ph.D. advisor, Marcel B\"okstedt, for both
suggesting this project and helping me with plenty of meetings and long
discussions when I was stuck, as well as providing
encouragement and good advice when it was needed most.
I would also like to thank my office mates and fellow Ph.D. students
at the Department of Mathematics
for both providing valuable feedback on ideas, and for keeping me
motivated.
Finally, I would like to thank all of my friends and family for their
support and friendship, without which my Ph.D. would have been much
more of a chore to write.

\section{Notation and conventions}

Throughout the dissertation, the following conventions will be used.
For two topological spaces $X$ and $Y$, we will write $X \cong Y$ to
denote that $X$ and $Y$ are homeomorphic, and $X\simeq Y$ to denote
that they are homotopy equivalent. Likewise, $f \simeq g$ means that
the maps $f$ and $g$ are homotopic. For subsets, we will use $A
\subset X$ when $A$ is any subset of $X$, and we will generally not
make a distinction between points of $X$ and one-point subsets of
$X$. The set difference of $X$ and $A$ will be written as $X-A$.

We will be working almost exclusively with topology in the complex
numbers,
so the disc $D\subset\C$ is the unit disc in the complex plane and
$D_p$ is the unit disc in $\C^p$. The non-zero elements in $\C$ will
be written as $\C^* = \C - 0$. $\GL_m(\C)$ will denote the group of
linear automorphisms of $\C^m$ and will be written as matrices with
respect to the standard basis $e_1,\dots,e_m$ of $\C^m$,
defined by
\[ \xy (9,0.5)*{\scriptstyle\leftarrow i};
(0,0)*{e_i =
\begin{pmatrix}
  0 \\
  \vdots \\
  0 \\
  1 \\
  0 \\
  \vdots \\
  0
\end{pmatrix}}; \endxy
%\big(0, \dots, 0, \overset{\substack{i\\\downarrow}}{1}, 0\dots,
%0\big).
\]

Since we are mainly concerned with algebraic topology, a space will
occasionally be replaced with a homotopy equivalent one without
mention. For example, $\T^n$ will denote both the complex units
$(\C^*)^n$, the
diagonal matrices in $\GL_n(\C)$, and $(S^1)^n$ depending
on context, and
the
punctured disc $D_p-0$ will sometimes be identified with the sphere
$S^{2p-1} \subset \C^p$. Likewise, we will not distinguish between the
different models of the flag manifold of $\C^m$, given as $\GLB{m}
\cong \UT{m} \cong \SUT{m}$ with $\B_m \subset \GL_m(\C)$ the upper (or
lower) triangular matrices in $\GL_m(\C)$, $\T^m \subset \U_m$ the
diagonal
matrices in the unitary group and $\T^{m-1}\subset \SU_m$ the diagonal
matrices in the special unitary group.

Homology and cohomology groups will always be singular groups
with integer coefficients. If a theorem is mentioned by name without a
reference, it can be found in \cite{hatcher}.

We will use permutations a great deal. The permutations of the set
$\set{1,\dots,m}$ will be denoted $S_m$.
These will be written as products of cycles, where we multiply by
considering permutations as functions and then composing. For
example, the
permutation $\sigma\in S_3$ given by
\begin{align*}
  \sigma(1) &= 3, \\
  \sigma(2) &= 1, \\
  \sigma(3) &= 2,
\end{align*}
will be written as $\sigma = (1 \; 3 \; 2) = (2\; 3)\cdot (1\; 2)$.
For our purposes, it will be useful to occasionally consider
a permutation as a matrix. We will do this by working with the
standard
basis, $e_1,\dots,e_m\in \C^m$. An element $\sigma \in S_m$ then
corresponds to the matrix defined by
\[ \sigma \cdot e_i = e_{\sigma(i)}. \]
So a permutation multiplied on a matrix from the left permutes the
rows of the matrix.


\vspace{5em}
\begin{flushright}
  Sarah Stoltze \\
  Aarhus University\\
  October, 2016
\end{flushright}




%%% Local Variables:
%%% mode: latex
%%% TeX-master: "main"
%%% End:
 % Dette er en afhandling baseret på mit arbejde som
                  % studerende blablabla. Tak til alt og alle (Marcel,
                  % muligvis institutet for støtte, proofreading,
                  % yadda) omkring mig for at gøre det muligt, etc.


\chapter{Résumé} 

I denne tekst behandler vi rum $Y_{m,n}(\sigma)\subset (\C^m)^n$,
defineret ud fra
en permutation $\sigma\in S_m$, og
kvotienter af disse,
\[ X_{m,n}(\sigma) = Y_{m,n}(\sigma) /\T^n. \]
Disse rum har en masse pæne egenskaber, og vi bruger en del tid på at
undersøge deres symmetrier, at studere deres topologi, og se på hvordan
man kan bevæge sig mellem
rummene. Specielt viser det sig at man kan danne $X_{m,n+1}(\sigma)$ ud
fra $X_{m,n}(\tau)$ ved at betragte de rigtige permutationer
$\tau\in S_m$. Denne observation fører til to vigtige værktøjer til at
studere rummene med. For det første giver det en spektralfølge der gør
at vi kan udregne
kohomologi-grupperne for $X_{m,n+1}(\sigma)$ hvis vi kender dem for
rummene $X_{m,n}(\tau)$, hvilket vi bruger til at udregne alle
kohomologi-grupperne for rummet $X_{3,3}$. For det andet får vi en
stabiliserings-afbildning
\[ s : X_{m,n} \to X_{m,n+1} \]
der gør at vi kan tage en grænseværdi af rummene og ende med det
uendelige rum $X_{m,\infty}$. Dette rum viser vi er ækvivalent til
løkkerummet for den fuldstændige flag-mangfoldighed $\SUT{m}$, hvilket
giver os en måde at studere løkkerummet ud fra de approksimerende rum
$X_{m,n}$, men også tillader os at gå den anden vej og få information
om $X_{m,n}$ ud fra løkkerummet. Denne ækvivalens bliver især
studeret for $m = 3$, og specielt er et af hovedresultaterne i
kapitel \ref{chap:loekker} at vi for $n \geq k$ har en isomorfi mellem
kohomologi-grupperne $H^k(X_{3,n})$ og kohomologi-grupperne 
$H^k(\Omega(\SU_3/\T^2))$.

En mere detaljeret gennemgang af indholdet i afhandlingen kan
findes i den engelske introduktion.

\chapter{Abstract}

This text treats spaces $Y_{m,n}(\sigma) \subset (\C^m)^n$, defined by
a permutation $\sigma\in S_m$, and quotients of these spaces,
\[ X_{m,n}(\sigma) = Y_{m,n}(\sigma) / \T^n. \]
These spaces have a lot of very nice properties, and we will spend
some time investigating various symmetries, studying their topology,
and looking at how you can pass from one space to another. In
particular, it turns out that you can build $X_{m,n+1}(\sigma)$ from
the spaces $X_{m,n}(\tau)$ by looking at the right permutations $\tau
\in S_m$. This leads to two important tools that can be used to study
the spaces. Firstly, it gives us a spectral sequence that allows us to
calculate the cohomology of $X_{m,n+1}(\sigma)$ if we know the
cohomology of $X_{m,n}(\tau)$, which we apply to the space $X_{3,3}$
to give a full description of the cohomology. Secondly, we get a
stabilisation map
\[ s : X_{m,n} \to X_{m,n+1} \]
that allows us to take a limit and end up with the infinite version of
the spaces, $X_{m,\infty}$. We show that this space is equivalent to
the loop space of the complete flag manifold $\SUT{m}$, which allows
us to study the loop space by looking at the simpler approximating
spaces $X_{m,n}$, but also allows us to get information about
$X_{m,n}$ by studying the loop space. This equivalence is studied
for $m=3$, and one of the main results of Chapter \ref{chap:loekker}
is that for $n \geq k$ we get an isomorphism between the cohomology
groups $H^k(X_{3,n})$ and the cohomology groups
$H^k(\Omega(\SU_3/\T^2))$.

A more detailed description of the contents of the dissertation can be
found in the introduction.

%%%%%% Previous version, can probably be reused.

% The following document is part of the midway exam for Ph.D. students
% at Aarhus University. It describes the work I have done on my project
% so far, and
% ends with a small presentation of ideas for future work. Due to the
% limited number of pages, the report is meant more as an overview and
% will occasionally only sketch a proof or argument instead of giving
% all the details.

% I would like to thank Thomas Schmidt for proofreading an earlier
% version of this report, and my advisor Marcel B\"okstedt for his help
% and inspiration during the past two years.

% \subsubsection{Conventions and notation}

% What follows is a brief description of the notation used in the
% report.
% For two topological spaces $X$ and $Y$, we will write $X \cong Y$ to
% denote that $X$ and $Y$ are homeomorphic and $X\simeq Y$ to denote
% that they are homotopy equivalent. Likewise, $f \simeq g$ means that
% the maps $f$ and $g$ are homotopic. For subsets, we will use $A
% \subset X$ when $A$ is any subset of $X$, and we will generally not
% make a distinction between points of $X$ and one-point subsets of
% $X$. The set difference of $X$ and $A$ will be written as $X-A$. With
% a distinguished basepoint $x_0$ in $X$, the set of based loops in $X$
% will be denoted $\Omega(X) = \set{ \g : S^1 \to X \mid \g(1) = x_0}$,
% and this comes equipped with the compact-open topology. The one-point
% compactification of $X$ is $X^+ = X\cup \set{\infty}$.

% We will be working exclusively with topology in the complex numbers,
% so the disc $D$ is the unit disc in the complex plane $\C$ and $D_p$
% is the unit disc in $\C^p$. The non-zero elements in $\C$ will be
% written as $\C^* = \C - 0$. $\GL_m$ will denote the group of
% linear automorphisms of $\C^m$ and the standard basis of $\C^m$ is
% $e_1,\dots,e_m$. Linear maps will be written as matrices with respect
% to the standard basis.

% Since we are mainly concerned with algebraic topology, a space will
% occasionally be replaced with a homotopy equivalent one without
% mention. For example, $T^n$ will denote both the torus $(\C^*)^n
% \subset \GL_n$ and $(S^1)^n\subset \C^n$ depending on context, and the
% punctured disc $D_p-0$ will sometimes be identified with the sphere
% $S^{2p-1} \subset \C^p$. Likewise, we will not distinguish between the
% different models of the flag manifold of $\C^m$, given as $\GLB{m}
% \cong \UT{m} \cong \SUT{m}$ with $B_m \subset \GL_m$ the upper (or
% lower) triangular matrices in $\GL_m$, $T^m \subset \U_m$ the diagonal
% matrices in the unitary group and $T^{m-1}\subset \SU_m$ the diagonal
% matrices in the special unitary group.

% Homology and cohomology groups will always be singular groups
% with integer coefficients. If a theorem is mentioned by name without a
% reference, it can be found in \cite{hatcher}.

% %\vfill % Er allerede så langt nede som muligt.
% \begin{flushright}
%   Simon Stoltze \\
%   Aarhus University\\
%   January, 2015
% \end{flushright}



%%% Local Variables: 
%%% mode: latex
%%% TeX-master: "main"
%%% End: 
 % Søren og JK har danske abstracts/resumé'er 
                   % der forklarer ca. det samme som deres
                   % introduktioner, bare på dansk. Kan være at det er
                   % et krav. Det består nok ca. i at skrive det på
                   % engelsk og oversætte det bagefter, måske korte
                   % lidt ned på det. Det skal i hvert fald skrives
                   % om. 

\chapter{Introduction}
\label{chap:introduktion}

The spaces we will be working with were originally introduced
in a very different form, which will briefly be described here. Though
this is not used in the dissertation, it hints at a connection with
algebraic geometry.
Consider the Grassmann-manifold, consisting of $m$-dimensional
subspaces of $\C^{\gdim}$:
\[ \mathrm{Gr}_m(\C^{\gdim}) = \set{V \subset \C^{\gdim} \mid
  \dim_\C(V) = m}. \]
This has a natural map to the projective space
\[ P = \mathbb{P}(\Lambda^m(\C^{m+n})), \]
given by the Pl\"ucker embedding
\begin{align*}
  \imath : \Grass &\to P \\
  \spa(v_1,v_2,\dots,v_m) &\mapsto v_1\wedge v_2\wedge \dots
                            \wedge v_m.
\end{align*}

Over $P$ we have the tautological bundle $H^*\to P$, consisting of
pairs $(L,v)$ with $L \in P$ a line in $\Lambda^m(\C^{m+n})$ and $v\in
L$ a vector on the line. The dual of
this bundle, $H$, consists of linear maps from $H^*$ to $\C$, and 
this has global sections given by picking out one of the
coordinates of $v$. In our case, a coordinate is given by choosing $m$
numbers
\[ (i_1 < i_2 < \dots < i_m) \subset \set{1,\dots,m+n}, \]
and considering the coefficient of
\[ e_{i_1}\wedge e_{i_2}\wedge \dots \wedge e_{i_m} \in
\Lambda^m(\C^{m+n}), \]
for the standard basis $e_1,\dots,e_{m+n}$ on $\C^{m+n}$. So each
set of indices $I$ defines a section
\[ \delta_I : P \to H. \]
We will single out the coordinates defined by the index sets
\[ (i < i+1 < \dots < i+m-1), \]
$i \in \set{1,\dots,m+n}$, where we replace $k$ by $k-m-n$ if
$k > m+n$. So we consider $m$ subseqeuent basis-vectors,
taken cyclically. Each of these coordinates defines a section
\[ \delta_i : P \to H, \]
so by tensoring them all together we get a section $\delta$ of the
tensor bundle $H^{\otimes(m+n)}$:
\[ \delta = \delta_1\otimes\delta_2\otimes\dots\otimes\delta_{m+n} : P
\to H^{\otimes(m+n)}. \]

In $P$ we have subspaces $P_I$, given by considering the lines where
the coordinate indexed by $I$ is zero. These are
projective spaces, of dimension one less than $P$, and they
exactly correspond to the inverse image $\delta_I^{-1}(0)$. We want to
work with the open set $U=P-\delta^{-1}(0)$ where our chosen
sections are all non-zero. If we use the Pl\"ucker embedding $\imath$,
we get a vector bundle
$E=\imath^{-1}(H^{\otimes(m+n)})$ on $\Grass$, a section
$\imath^{-1}(\delta) : \Grass \to E$, and an open set
$X=\imath^{-1}(U)$ in $\Grass$ that exactly corresponds to the points
on which the section is non-zero. By considering how the embedding
$\imath$ is defined, we can see that the elements of $X$ are
the subspaces that have a basis $v_1,\dots,v_m$ such that if we
consider
them as row vectors and calculate the determinants of $m$ subsequent
columns, repeated cyclically, we always get a non-zero number. In
symbols, this means
\[ D_i(v_1,\dots,v_m) = \det
 \begin{pmatrix}
   (v_1)_{i} & (v_1)_{i+1} & \dots  & (v_1)_{i+m-1} \\
   (v_2)_{i} & (v_2)_{i+1} & \dots  & (v_2)_{i+m-1} \\
   \vdots  & \vdots    & \ddots & \vdots       \\
   (v_m)_{i} & (v_m)_{i+1} & \dots  & (v_m)_{i+m-1}
 \end{pmatrix} \neq 0 \]
for all $i \in \set{1,\dots,m+n}$. Since $D_1$ is non-zero, we can
can choose a basis where the first $m$
columns are the identity matrix without affecting the span. If we
want to preserve this property, we
can only change the basis by scaling each column by a scalar. So our
space $X$ can be represented by sequences of $n$ vectors,
\[ v_1,\dots,v_n \in \C^m, \]
such that if
we insert the identity matrix before and after the sequence,
\[ e_1,\dots,e_m,v_1,\dots,v_n,e_1,\dots,e_m, \]
any $m$ subsequent vectors in this new sequence are linearly
independent, and we have quotiented out by an action of the torus
$\T^n$ given by letting the $i$'th coordinate of the torus act on
$v_i$. This is the space $X_{m,n}$ described in Definition
\ref{def:kvotientrum}.

The spaces
were originally studied in a slightly simpler form by my advisor,
Marcel B\"okstedt. He looked at sequences of vectors $(v_1,\dots,v_n)$
in $\C^2$, with the requirement that $v_i$ and $v_{i+1}$ are linearly
independent, that the first entry of $v_1$  is non-zero, and either
the first or the second entry of $v_n$ is non-zero as well. Using the
notation introduced in this thesis, these are the spaces $Y_{2,n}$ or
$Y_{2,n}((1\;2))$ from Definition \ref{def:rum}. Taking a quotient
of the action
\begin{align*}
  \T^n \times Y_{2,n}(\sigma) &\to Y_{2,n}(\sigma) \\
  (\l_1,\dots,\l_n)\cdot (v_1,\dots,v_n) &= (\l_1v_1,\dots,\l_nv_n),
\end{align*}
we end up with the spaces referred to as $X_{2,n}(\sigma)$ in
Definition \ref{def:kvotientrum}. He then showed that $Y_{2,n}$ was
covered by the spaces $Y_{2,n-1}\times\T^2$ and $Y_{2,n-1}((1\; 2))\times
\T$, giving inclusion maps
\[ s : Y_{2,n-1} \cong Y_{2,n-1}\times\set{(1,1)} \to Y_{2,n}. \]
By using the Mayer-Vietoris sequence for the cover, and the
corresponding cover of $X_{2,n}$, he proved the
following theorem.
\begin{theorem}
  The space $X_{2,n}$ has homology groups
  \[ H_q(X_{2,n}) \cong
  \begin{cases}
    \Z & 0 \leq q \leq n, \\
    0 & \text{otherwise.}
  \end{cases} \]
  The inclusion
  \[ s : X_{2,n-1} \to X_{2,n} \]
  is an isomorphism on homology when the degree is $\leq n-1$.
\end{theorem}

The homology of $X_{2,n}$ consists of a copy of the integers in all
degrees up to $n$, so when $n$ is large it looks like the homology
groups of the loop space
$\Omega(\SU_2/\T)$, which has a copy of the integers in each
degree. We can define a map from $X_{2,n}$ to the loop space
by sending $(v_1,\dots,v_n)$ to
the path given by connecting up the matrices
\[ \Id \leadsto (e_2,v_1) \leadsto (v_1,v_2) \leadsto \dots \leadsto
\Id, \]
and it turns out that this map is an isomorphism on homology when the
groups are non-zero.

The aim of this dissertation is to extend some of
these results to $X_{m,n}$, $m > 2$, with a focus on cohomology
instead of homology. The cohomology groups become much harder to
calculate when $m > 2$, since the Mayer-Vietoris sequence that was
used for $X_{2,n}$
does not give enough information. Instead we have to cover $X_{m,n}$
with several subspaces, which means we have to use spectral
sequences to compute the cohomology groups and this leads to increased
complexity. However, it is still possible to prove that $X_{m,n}$
approximates the loop space $\Omega(\SUT{m})$ when $n$ is large, which
is the content of Theorem \ref{thm:loekker}. This approximation allows
us to compute $H^q(X_{3,n})$ when $n \geq q$, as is done in Theorem
\ref{thm:kostab}.

The thesis starts with a short introduction to some of the tools and
subjects that
will be required in the following chapters, contained in Chapter
\ref{chap:generelt}. Each section consists of a brief description of
a subject, with references to
important theorems needed for our purposes. The topics that are
covered are spectral
sequences, loop spaces, flags, and Coxeter groups, and since it is not
possible to give a full overview of any of these in a couple of
pages, every section also contain references to literature with a
more extensive coverage.

In Chapter \ref{chap:rum}, we give a description of the spaces we will
be working with, first introducing subspaces $Y_{m,n}(\sigma)$ of
$(\C^{m})^{n}$, which depend on permutations $\sigma\in S_m$, and
then quotienting out with an action of the complex torus $\T^n =
(\C^*)^n$, giving the spaces $X_{m,n}(\sigma)$ that we will mostly be
working with. We show that these spaces are non-empty when $n$ is
large enough, and we spend some time working out how the space
$Y_{m,n}(\sigma)$ can be built from the spaces $Y_{m,n-1}(\tau)$, and
exactly which permutations $\tau$ are necessary to do this in Theorem
\ref{thm:permutation}. This
allows us to take a limit, and gives us the spaces
$Y_{m,\infty}(\sigma)$ and $X_{m,\infty}(\sigma)$. We end
with a very brief discussion of the spaces we get when we replace the
complex numbers with the reals or quaternions.

Chapter \ref{chap:udregninger} starts with a description of a spectral
sequence that will be used extensively. This allows us to compute the
cohomology of the space $Y_{m,n}(\sigma)$ from the cohomology of the
spaces $Y_{m,n-1}(\tau)$ that it is built from. We immediately use
this spectral sequence to work out the first cohomology group of both
$Y_{m,n}$ and $X_{m,n}$ for all $m$ and $n$. The rest of the
chapter is spent doing computations with the spectral sequence for the
spaces $Y_{3,3}$ and
$X_{3,3}$, ending with a complete description of the cohomology
groups for these spaces.

The following chapter, Chapter \ref{chap:loekker}, studies the limit
spaces $X_{m,\infty}(\sigma)$ in detail. We start by showing that they
are homotopy equivalent to the loop space $\Omega(\SUT{m})$, which
gives us the cohomology of the limit space. The proof does not depend
on working over the complex numbers, so it easily carries over to the
real and quaternionic cases. We then give a brief
mention of other ways we could define the limit space, which are all
equivalent from a homotopy theoretic standpoint. By working
with one of these other stabilisations of $X_{3,n}(\sigma)$, we get
Theorem \ref{thm:kostab}
that tells us when the cohomology groups of
$X_{3,n}(\sigma)$ stabilise to the cohomology groups of
$X_{3,\infty}(\sigma)$.

The thesis concludes with a brief description of the open questions
still remaining, contained in Chapter \ref{chap:fremtid}. These are
questions that have been raised during the last parts of the project
and thus have not been studied in detail, as well as small ideas for
how the results could be applied to the free loop space instead of
the space of based loops considered in the thesis.



% En introduktion til afhandlingen, inklusive kapitel-oversigt og
% ideer bag. Her bør for eksempel indgå en forklaring af det der pjat
% med vektorbundter som Marcel har forsøgt at forklare før.
% Derudover nok også en oversigt over de vigtigste resultater der
% indgår.

%%%%% Previous
% \chapter{Introduction}

% % Noget om starten, hvorfor dette er interessant. Se bagside af blok
% % Noget om notation... eg \subset, D-0, hvad er D_p
% % Alting komplekst, så \GL_m giver mening etc. Måske også de
% % forskellige versioner af flag-mangfoldigheder.
% % Og noget med at T=\C^*....

% This document gives a description of a family of spaces $\X\subset
% \C^{mn}$, indexed by two natural numbers $m$ and $n$, and a quotient
% of these spaces, $Y_{m,n} = \X / T^n$. The spaces were
% originally defined as the zero-set of a section on a vector bundle of
% a Grassmann manifold, and were studied to understand the properties of
% families of Calabi-Yau manifolds. However, the spaces turned out to
% have other properties that made them interesting in their own
% right. The most important property is that they can be used to form a
% space that is homotopy equivalent to a space that appears in physics,
% namely the loop space $\Omega(\SUT{m})$. This hints at a way of
% extracting information about the loop space by working with the
% simpler spaces $Y_{m,n}$.

% The spaces have been studied by Marcel B\"okstedt, and some of his
% results are gathered in \cite{bokstedt}. In particular it includes the
% following theorem that covers the special case where $m=2$:
% \begin{theorem}
%   The homology of $Y_{2,n}$ is
%   \[ H_q(Y_{2,n}) =
%   \begin{cases}
%     \Z & 0 \leq q \leq n \\
%     0 & \text{otherwise}
%   \end{cases} \]
%   There is a stabilization map $s : Y_{2,n} \to Y_{2,n+1}$ that
%   induces an isomorphism
%   \[ s_* : H_q(Y_{2,n}) \to H_q(Y_{2,n+1}) \]
%   for $q \leq n$. There is also a map $f : Y_{2,n} \to
%   \Omega(\mathrm{SU_2}/T^1)$ that induces an isomorphism in homology in
%   degree $q\leq n$.
% \end{theorem}

% The aim of the following is to try and extend these results to the
% general case of $Y_{m,n}$, with $m\geq 2$.

% This progress report is divided into three main parts.
% The first chapter gives an introduction to the spaces we are
% interested in and lists some of the properties that will be used
% in the other chapters.
% The second chapter is mainly concerned with an explicit calculation of
% the cohomology of the spaces $X_{3,3}$ and $Y_{3,3}$. This is done by
% applying the spectral sequence coming from a filtration of $X_{3,3}$,
% computing the various relative cohomology groups needed, and
% finding the differentials of the sequence.
% The third and final chapter concerns the relation to the loop space
% described above. Here it is shown that by taking a limit of the spaces
% $Y_{m,n}$, we get a space homotopy equivalent to the loop space. The
% chapter concludes with a brief overview of plans for future studies.

%%% Local Variables: 
%%% mode: latex
%%% TeX-master: "main"
%%% End: 
 % Nu på engelsk. En generel introduktion til
                       % ideerne bag og hvad vi har forsøgt at
                       % vise. 2-3 sider lader til at være det
                       % normale. Indholder også typisk en oversigt
                       % over de resterende kapitler og lidt
                       % forklaring af hvad de går ud på.
\mainmatter

% Et kapitel med lidt generelle resultater. Dette vil nok mest bestå
% af lidt om spektralfølger (Hatcher?) og lidt om løkkerum (Milnor).

% Her skal samles lidt generel information om SS og løkkerum

\chapter{Spectral sequences and loop spaces}
\fixme{Bedre titel...}

The purpose of this chapter is to give a brief introduction to two
important topics in this thesis, namely spectral sequences and loop
spaces. Spectral sequences are an important tool in algebraic topology
and can be considered a generalisation of long exact sequences, while
loop spaces are a class of spaces that appear naturally when
considering the spaces that will be studied in this thesis. The
following pages will collect some basic info and state the results
that are of interest in next chapters, while providing references for
proofs and further information. \fixme{Måske noget med at vi starter
  ca. fra Hatcher, alt andet betragtes som ``nyt''}

\section{Spectral sequences}
\label{sec:ss}

A spectral sequence, as we will be using them, is a tool for computing
the cohomology groups of a topological space, $X$, by considering an
increasing sequence of subspaces
\[ X_0 \subset X_1 \subset X_2 \subset \dots \subset X, \]
with the union covering all of $X$,
\[ X = \bigcup_{i} X_i. \]
An example that most people should be familiar with is the long exact
sequence of a pair, where
we only have a single subspace $A \subset X$. Another simple example
is the case where there are two subspaces,
\[ X = A \cup B, \]
in which case the cohomology of $X$ can often be computed from the
cohomology of $A$, $B$ and $A\cap B$ by using the Mayer-Vietoris
sequence. This generalises to a spectral sequence to compute
the cohomology of $X$ by working with an abitrary covering. As might
be expected, since spectral sequences work more generally they also
carry a corresponding increase in complexity when using them for
calculations. In the following section we will focus on the specific
case that 
we will be using in this text. For a general introduction to the
subject, see either \cite{hatcherss} or \cite{mccleary}.

\subsection{The spectral sequence of a filtration}
\label{sec:ss-filtration}

Consider a space $X$ with a filtration, that is an increasing sequence
of subspaces that cover $X$:
\[ X_0 \subset X_1 \subset X_2 \subset
\dots \subset X_k  = X. \]
All of the following can be modified to work if the filtration is not
finite, but we will not need this here. We will augment the filtration
by defining additional subspaces,
\[ X_s =
\begin{cases}
  \emptyset & s \leq -1 \\
  X & s \geq k+1
\end{cases}. \]

With such a sequence, we can
consider the relative cohomology groups
\[ E_1^{p,q} = H^{p+q}(X_p,X_{p-1}), \]
and the differentials between them coming from the long exact sequence
of the triple $(X_{p+1},X_p,X_{p-1})$:
\[ d_1 : E_1^{p,q} = H^{p+q}(X_p,X_{p-1}) \to H^{p+q+1}(X_{p+1},X_p) =
E_1^{p+1,q}. \]
This map factors as the composition of an
inclusion of the pair $(X_p,X_{p-1})$ and the boundary map of the pair
$(X_{p+1},X_p)$,
\[ \xymatrix{ d_1 : H^{p+q}(X_p,X_{p-1}) \ar[r]^<<<<<{\jmath^*} &
  H^{p+q}(X_p) \ar[r]^<<<<<{\delta} & H^{p+q+1}(X_{p+1},X_p),
} \]
and in particular we get that the composition of two such maps must be
zero since the composition $\delta \circ \jmath^* = 0$ by
exactness. So we can compute the homology of these groups and get new
groups,
\[ E_2^{p,q} = \ker d_1 / \im d_1. \]
More suprisingly, we can also find new maps
\[ d_2 : E_2^{p,q} \to E_2^{p+2,q-1} \]
that make these new groups into a chain complex, allowing us to form
\[ E_3^{p,q} = \ker d_2 / \im d_2. \]
This process can be continued, giving a sequence of groups and maps
\[ \xymatrix{d_r : E_r^{p,q} \ar[r] & E_r^{p+r,q-r+1} } \]
where the composition $d_r \circ d_r$ is zero and
$E_r^{p,q}$ is formed by taking the homology of $E_{r-1}^{p,q}$ with
the map $d_{r-1}$. Since we have $E_r^{p,q} = 0$ for $p < 0$ or $p >
k$, these maps will eventually always go either to or from a zero
group and hence the groups $E_r^{p,q}$ will not change when we
increase $r$. This allows us to
define the stable groups $E_\infty^{p,q}$ as these $r$-independent
groups. The following theorem relates them to the cohomology groups of
$X$ that we are interested in calculating.

\begin{theorem}
  \label{thm:ss}
  Given a filtration $\set{X_s}$ of a topological space $X$ as above,
  there is a sequence of abelian groups $E_r^{p,q}$ and maps
  \[ d_r : E_r^{p,q} \to E_r^{p+r,q-r+1} \]
  that satisfy the following:
  \begin{itemize}
  \item The first groups are
    \[ E_1^{p,q} = H^{p+q}(X_p,X_{p-1}) \]
    and the map $d_1$ is the boundary map of the triples
    $(X_{p+1},X_p,X_{p-1})$.
  \item The composition $d_r \circ d_r$ is zero and
    \[ E_{r+1}^{p,q} = \ker d_r / \im d_r, \]
    computed at $E_r^{p,q}$.
  \item There is a filtration of the $n\!$'th cohomology group of $X$,
    \[ 0 \subset F_n^n \subset \dots F_0^n = H^n(X), \]
    with the limit groups $E_\infty^{p,n-p}$ isomorphic to the
    quotients
    \[ E_\infty^{p,n-p} \cong F_p^n / F_{p+1}^n. \]
  \end{itemize}
\end{theorem}

This theorem is a combination of Proposition 1.2 and Theorem 1.14 in
\cite{hatcherss}, and the proof can be found there as well. Assuming
we can compute the spectral sequence, the group $H^n(X)$ can then be
found by solving extension problems of the form
\[ 0 \to F_{p+1}^n \to F_p^n \to E_\infty^{p,n-p} \to 0, \]
starting from the short exact sequence
\[ 0 \to E_\infty^{n,0} \to F_{n-1}^n \to E_\infty^{n-1,1} \to 0. \]
For example, if we have that the groups $E_{\infty}^{p,n-p}$ are free
for all $p$, the short exact sequences always split and $H^n(X)$ is
isomorphic to the direct sum of these groups,
\[ H^n(X) \cong \bigoplus_p E_\infty^{p,n-p}. \]

\section{Loop spaces}
\label{sec:ls}

Throughout this section, every topological space will have a
distinguished basepoint. We will be looking a spaces of loops,
consiting of maps from the circle to a given space $X$, preserving the
basepoint. For a general introduction to loop spaces, consult e.g.
\cite{may}. The results in this section are taken from
\cite{milnor} and the proofs can be found there as well.

\begin{definition}
  If $X$ is a topological space with basepoint $*\in
  X$, the loop space $\Omega(X)$ is the space of based maps from the
  cicle to $X$,
  \[ \Omega(X) = \set{ \g : (S^1,1)\to (X,*) } = \set{ \g : [0,1] \to X
    \mid f(0) = f(1) = * }. \]

  We equip this set with the compact-open topology, so for a compact
  set $K \subset S^1$ and an open set $U \subset X$, we get an open
  subset of the loop space consisting of functions mapping $K$ to $U$,
  \[ V(K,U) = \set{ \g\in \Omega(X) \mid \g(K) \subset U }. \]
  The collection of all such sets form a subbase for the compact-open
  topology on $\Omega(X)$.
\end{definition}

We will only worry about the case where $X$ is a Riemannian manifold,
so it comes equipped with a Riemannian metric $g$ and an induced
topological metric $d$. In this case, the compact-open topology is
induced by the metric
\[ d^*(\g,\g') = \max_t d(\g(t),\g'(t)). \]
By \cite[Theorem 17.1]{milnor}, this is homotopy equivalent to the
space of piecewise smooth loops in $X$, and we will generally work
with this space instead. The proof will not be repeated here, but
there is a construction in the proof that deserves a mention. For an
open cover $\set{X_\alpha}_{\a\in I}$ of $X$, we define subsets of the
loop space by requiring that the loops are piecewise contained in a
set in the cover,
\[ \Omega_k(X) = \set{\g \in \Omega(X) \delim \Big\vert\delim
  \text{For each } j \leq 2^k, \g_{\left|\left[
        \frac{j-1}{2^k},\frac{j}{2^k} \right]\right.} \text{ is 
    contained in $X_\alpha$ for some $\alpha$.} }. \]
This gives an increasing sequence of open subsets of $\Omega(X)$
that cover the entire space. In particular, a compact subset of
$\Omega(X)$ will be contained in $\Omega_k(X)$ for some $k$. By
defining a cover of $X$ with good properties, we will be able to work
exclusively with loops where we have some control
over their behaviour.


% Nævne at vi ikke behøver køre * -> *, men også kan køre p -> q

Note that instead of considering loops, we could also consider the
space of paths between two points $p,q\in X$, given by
\[ \Omega(X,p,q) = \set{ \g : [0,1] \to X \mid \g(0) = p, \g(1) = q
}. \]
But by composing such a path with a fixed path $\a$ from $q$ to $p$,
we end up with a loop at $p$. If we compose a loop at $p$ with the
path from $p$ to $q$ given by running $\a$ in reverse, we get a path
from $p$ to $q$, and the composition of these two operrations is
homotopic to the identity. This shows that all such spaces 
are homotopy equivalent, so we will not make a major distinction
between them.

\fixme{Et kort indblik i flag, ca. en halv eller en hel side?}

\section{Flags}
\label{sec:flags}

The last thing to be discussed is the concept of a flag in a finite
dimensional vector space. In an
$m$-dimensional vector space $V$, a flag is an increasing sequence of
subspaces of $V$,
\[ V_1 \subset V_2 \subset \dots \subset V_m = V, \]
with the dimension of $V_i$ equal to $i$. So the dimension of each
subspace is one greater than the last. Note that if we pick a non-zero
vector $v_1$ in $V_1$, another (non-zero) vector $v_2$ in $V_2 - V_1$,
and continue picking $v_i$ in $V_i - V_{i-1}$, we end up with $m$
linearly independent vectors. This allows us to construct a matrix in
$\GL(V)$ by considering
\[ [v_1,\dots,v_m], \]
and we can recover the flag from the matrix by considering
\begin{align*}
  V_1 =\, &\spa(v_1) \\
  V_2 =\, &\spa(v_1,v_2) \\
  \vdots\,\, & \\
  V_m =\, &\spa(v_1,\dots,v_m).
\end{align*}
Likewise, we could take the span of the columns starting from $v_m$
and working backwards to get another flag.
\begin{definition}
  Given an invertible matrix $A = [v_1,\dots,v_m]$, the \textit{left
    flag} of $A$ is
  \begin{alignat*}{5}
    \mathrm{Fl_L}(A)\, &&= \Big(&\spa(v_1) &&\subset \spa(v_1,v_2)
    &&\subset \dots &&\subset \spa(v_1,\dots,v_m) = V\Big).
  \end{alignat*}
  The \textit{right flag} of $V$ is defined similarly as
  \begin{alignat*}{5}
    \mathrm{Fl_R}(A)\, &&= \Big(&\spa(v_m) &&\subset
    \spa(x_{v-1},x_v) &&\subset \dots &&\subset
    \spa(v_1,\dots,v_m) = V\Big). 
  \end{alignat*}
\end{definition}
Note that we could multiply any column of $A$ by a non-zero scalar and
we would end up with the same flag. Likewise, if we replace $v_j$ by
$v_j + \lambda v_i$ for a scalar $\lambda$ and indices $j > i$, we do not
change the left flag of $A$. The right flag is unchanged by doing
the same operation with $j < i$. In particular, if $L$ is a
lower-triangular invertible matrix we can do the above operations one
column at a time, starting from the last, and see that
\[ \mathrm{Fl_R}(L) = \mathrm{Fl_R}(\Id). \]
And similarly, if $U$ is upper-triangular and invertible and $\sigma$
is a permutation matrix:
\[ \mathrm{Fl_L}(\sigma U) = \mathrm{Fl_L}(\sigma). \]
This is slightly harder to see, but $\sigma$ permutes the rows of $U$,
moving the first row to row $\sigma(1)$, the second row to row
$\sigma(2)$, and so on. By adding the first column to remaining ones,
with a suitable scalar multiplied on, we can ensure that the only
non-zero entry in row $\sigma(1)$ of $\sigma U$ is in the first
column, row $\sigma(1)$, just as it is in the matrix
$\sigma$. Continuing like this gives the desired identity.
\fixme{Mere her... Eller en afslutning}

\section{Coxeter groups and the Bruhat order}

\fixme{Skriv lidt her... Bare lige nok til at vise at permutationerne
  er relaterede}

%%% Local Variables: 
%%% mode: latex
%%% TeX-master: "main"
%%% End: 


% Muligvis bør der også være et kapitel med en introduktion til
% spektral-følger og andet generelt brugt teori, såsom lidt
% Morse-teori fra Milnor, en introduktion til løkke-rum, måske lidt om
% Coxeter-grupper? Det kan diskuteres om det skal i samme kapitel som
% introduktionen til rummene.
% Her inkluderer de gamle mennesker små beviser men lader store
% sætninger stå urørte. Sikkert henvisen til sætninger man skal bruge
% direkte og så vise nogle af de specialtilfælde der dukker op. Eller
% nogle gange at give ideerne i de lange beviser.





% Skal der laves et intro-kapitel med baggrundsmateriale? Løkker
% (Milnor), Bruhat, spektralfølger.
% Ellers skal det nok ind her et sted, men måske spredt lidt ud over
% kapitlet.

\chapter{Spaces of linearly independent vectors}
\fixme{Bedre titel?}
\fixme{Sæt memoir til at bruge den rigtige chapter-overskrift så det
  ikke ligner lort.}

% Introduktion til rum
%% Ukvotienterede (Y_{m,n}) og de rigtige (X_{m,n} = Y_{m,n}/T^m)
%%% Så introducer gruppevirkningen her. Og vis at Y \cong X x T
%%% Bør nok også nævne virkningen på rækkerne, da den spiller ind i
%%% spektralfølgen.
%% Andre permutationer.
%% Andre koefficienter: \R og \H. Disse dukker først op igen
%% senere, men det meste af teorien i dette afsnit burde være
%% uafhængigt af koefficienterne, op til at nogle skalarer flytter
%% lidt rundt. Specielt kan vi lige så godt smide dem ind her.

In this chapter we will introduce the spaces that will be studied in
the rest of the document, and collect various useful
results.\fixme{Bleh. Skriv mere}

\section{The spaces} \fixme{Bedre overskrift...}
\label{sec:rum}
Before we can define the spaces that we will be working with, it will
be helpful to define an approximation to them.

\begin{definition}
  For natural numbers $m$ and $n$, the space $Y_{m,n}\subset \C^{mn}$
  is
  \[ Y_{m,n} = \set{(a_1,\dots,a_n) \in (\C^m)^n \delim\Bigg\vert\delim
    \begin{matrix}
      \text{Any } m \text{ subsequent vectors in } \\
      (e_1,\dots,e_m,a_1,\dots,a_n,e_1,\dots,e_m) \\
      \text{ are linearly independent.}
    \end{matrix} } \]
  where $e_1,\dots,e_m \in \C^m$ denotes the standard basis
  vectors,\fixme{Det der i...}
  \[ e_i =
  \big(0,\dots,0,\overset{\substack{i\\\downarrow}}{1},0\dots,0\big).\]

\end{definition}
We will switch between considering $(a_1,\dots,a_n)\in Y_{m,n}$ as an
ordered sequences of linearly independent vectors and an $m\times n$
matrix $[a_1,\dots,a_n]$ as needed and without mention. We can replace
the standard basis that appears in the definition of $Y_{m,n}$ with
any collection of linearly independent vectors and get similarly
defined spaces: 
\begin{definition}
  For two invertible matrices $X$ and $Y$ in
  $\GL_m(\C)$, the space $Y_{m,n}(X,Y) \subset \C^{mn}$ is
  \[ Y_{m,n}(X,Y) = \set{(a_1,\dots,a_n) \in (\C^m)^n
    \delim\Bigg\vert\delim 
    \begin{matrix}
      \text{Any } m \text{ subsequent vectors in } \\
      (x_1,\dots,x_m,a_1,\dots,a_n,y_1,\dots,y_m) \\
      \text{ are linearly independent.}
    \end{matrix} } \]
  where $(x_1,\dots,x_m)$ and $(y_1,\dots,y_m)$ are the column vectors
  of $X$ and $Y$ respectively.
  
  There are two special cases that will be especially important.
  If both $X$ and $Y$ are the identity-matrix, we get the previously
  defined spaces,
  \[ Y_{m,n} = Y_{m,n}(\Id,\Id). \]
  When $X$ is the identity-matrix, we will usually omit it from the
  notation and instead write
  \[ Y_{m,n}(Y) = Y_{m,n}(\Id,Y). \]
\end{definition}

Note that $Y_{m,n}(X,Y)$ can also be defined as the
preimage of the open set $(\C^*)^{m+n-1}$ under the continuous map
\begin{align*}
  \Det_{X,Y} : (\C^m)^n \to \C&^{m+n-1} \\
  \Det_{X,Y}(a_1,\dots,a_n) = (\det&(x_2,\dots,x_m,a_1),\\
                               \det&(x_3,\dots,x_m,a_1,a_2), \\
                               &\vdots\\
                               \det&(a_n,y_1,\dots,y_{m-1})).
\end{align*}\fixme{Se om alignment kan blive bedre på pilen.}
This shows that it is an open subset of $(\C^m)^n$. Another important
feature of $Y_{m,n}(X,Y)$ is that it does not really depend on the
particular matrices $X = [x_1,\dots,x_m]$ and $Y = [y_1,\dots,y_m]$,
but\fixme{hbox for $X$}
only on the two flags in $\C^{m}$ 
defined from the columns: \fixme{Snak lidt om flag tidligere/i introen?}
\begin{alignat*}{5}
  \mathrm{Fl_R}(X)\, &&= \Big(&\spa(x_m) &&\subset \spa(x_{m-1},x_m) &&\subset
    \dots &&\subset \C^m\Big) \\
  \mathrm{Fl_L}(Y)\, &&= \Big(&\spa(y_1) &&\subset \spa(y_1,y_2) &&\subset
    \dots &&\subset \C^m\Big).
\end{alignat*}
This is because we are considering linear independence of the vectors
\[ (x_1,\dots,x_m,a_1,\dots,a_n,y_1,\dots,y_n), \]
and if $(a_1,\dots,a_{k})$ is a collection of linearly independent
vectors of $\C^m$, then the vectors
$(x_{k+1},\dots,x_m,a_1,\dots,a_{k})$ are
linearly independent if and only if $a_1$ is not in the subspace
spanned by $(x_{k+1},\dots,x_m)$, $a_1$ and $a_2$ are linearly
independent and not in the
subspace spanned by $(x_{k+2},\dots,x_m)$ and so forth. The same
argument can be applied to the columns of $Y$.

At first glance it looks like we are considering a huge collection of
spaces, but most of them are indistinguishable topologically, as the
following lemma shows.

\begin{lemma} \label{lem:rum-perm}
  For $X,Y\in \GL_m(\C)$, there exists a unique permutation $\sigma
  \in S_m$ \fixme{Introducer notation for $S_m$} such that the space
  $Y_{m,n}(X,Y)$ is homeomorphic to $Y_{m,n}(\sigma)$.
\end{lemma}

\begin{proof}
  Since $X$ is invertible it preserves linear independence of vectors,
  and we get a homeomorphism
  \begin{align*}
    Y_{m,n}(X,Y) &\to Y_{m,n}(\Id,X^{-1}Y) \\
    (a_1,\dots,a_m) &\mapsto (X^{-1}a_1,\dots,X^{-1}a_m).
  \end{align*}
  The matrix $X^{-1}Y$ is invertible, and by using the
  Bruhat-decomposition of $\GL_m(\C)$ we can write it uniquely as a
  product \fixme{Reference, enten til tidligere eller til bog}
  \[ X^{-1}Y = L\sigma U \]
  where $L$ is lower-triangular, $U$ is upper-triangular and $\sigma$
  is a permutation. Hence we get a homeomorphism,
  \[ Y_{m,n}(X,Y) \cong Y_{m,n}(L^{-1}, \sigma U). \]
  But by the remark above, the space $Y_{m,n}(L^{-1},\sigma U)$ only
  depends
  on the flags \fixme{Introducer flag. Måske flyt dette til en
    bemærkning efter definitionen af $Y_{m,n}$. Og vis at højre er
    uændret ved mult med nedre-triangulære og venstre er uændret ved
    mult med øvre-triangulære}
  \begin{align*}
    \mathrm{Fl_R}(L^{-1}) &= \mathrm{Fl_R}(\Id) \\
    \mathrm{Fl_L}(\sigma U) &= \mathrm{Fl_L}(\sigma).
  \end{align*}
  So $Y_{m,n}(L^{-1},\sigma U)$ is the same space as $Y_{m,n}(\sigma)$.
\end{proof}

Due to this, we will almost exclusively consider the spaces
$Y_{m,n}(\sigma)$ throughout the text. To get a feel for the spaces,
we will now work with some simple examples.

\begin{example}
  \label{ex:n=1}
  The space $Y_{m,1}$ is homeomorphic to $(\C^*)^m$, since
  \[\begin{pmatrix}
    1&0&\dots&0&\l_1&1&0&\dots&0 \\
    0&1&\dots&0&\l_2&0&1&\dots&0 \\
    \vdots&\vdots&\ddots&\vdots&\vdots&\vdots&\vdots&\ddots&\vdots \\
    0&0&\dots&1&\l_m&0&0&\dots&1
  \end{pmatrix} \]
  has linearly independent columns if and only if all the scalars
  $\l_1,\dots,\l_m$ are invertible.
\end{example}

\begin{example}
  The space $Y_{2,2}$ is
  \[ \set{
    \begin{pmatrix}
      a & b \\
      c & d
    \end{pmatrix} \delim\bigg\vert\delim a \neq 0, d \neq 0, ad-bc
    \neq 0 } \cong (\C^*)^2 \times \set{(b,c)\in\C^2 \delim\vert\delim
    bc \neq 1}, \]
  with the homeomorphism given by
  \[
  \begin{pmatrix}
    a & b\\
    c & d
  \end{pmatrix} \mapsto \left(a,d,\frac{b}{d},\frac{c}{a}\right).\]
  The product of the two fractions is
  \[ \frac{bc}{ad} = \frac{bc-ad}{ad} + 1 = 1-\frac{ad-bc}{ad} \neq
  1, \]
  since $ad-bc \neq 0$.
\end{example}

\begin{example}
  The space $Y_{3,2}$ is
  \[\set{
  \begin{pmatrix}
    a &b \\
    c & d \\
    e & f
  \end{pmatrix} \delim\Bigg\vert\delim a \neq 0, f \neq 0, ad-bc \neq
  0, cf-ed \neq 0}.
  \]
  This is homeomorphic to $(\C^*)^4\times \set{(x,y,z)\in\C^3 \mid
    y-z-xyz = 1}$, with the map given by
  \[
  \begin{pmatrix}
    a &b \\
    c & d \\
    e & f
  \end{pmatrix} \mapsto
  \left(a, f, D=\frac{ad-bc}{af}, C=\frac{cf-de}{af}, x=\frac{bC}{fD},
    y=\frac{c}{aC}, z=\frac{eD}{aC} \right). \]
  Some of the coordinates have been named to shorten the writing
  slightly. By direct calculation,
  \begin{align*}
    y-z-xyz &= \frac{c}{aC}-\frac{eD}{aC}-\frac{bce}{a^2fC} \\
            &= \frac{acf-aefD-bce}{a^2fC} \\
            &= \frac{acf-e(ad-bc)-bce}{a(cf-de)}  \\
            &= \frac{acf-ade}{acf-ade} \\
            &= 1.
  \end{align*}
  The inverse map is
  \[ (a,f,D,C,x,y,z) \mapsto
  \begin{pmatrix}
    a & \frac{fDx}{C} \\
    aCy & fD(1+xy) \\
    \frac{aCz}{D} & f
  \end{pmatrix}. \]
\end{example}


% Vis at de er ikke-tomme
%% Nok både dimensions-argumentet og vise at stabiliseringen af e_1 er
%% i alle rummene.

To work with the above spaces, it will help to know that they are not
empty. For example, by considering the permutation $\big(1\; 3\big)$,
we see that it is impossible to satisfy the 
conditions for being in $Y_{3,1}\left(\big(1\; 3\big)\right)$. To see
this, take any vector $a_1$ in $\C^3$ and look at
$(e_1,e_2,e_3,a_1,e_3,e_2,e_1)$. If this is to be in the space, it
must have any three subsequent vectors linearly independent, but
$(e_3,a_1,e_3)$ is not invertible no matter what $a_1$ is. This shows
that the choice of $X$ and $Y$ in $Y_{m,n}(X,Y)$ plays a role in
determining when the spaces are empty. But as we shall now see, the
spaces are non-empty as long as we are not in an obviously empty
space. The exact meaning of ``obviously'' is given in the following
lemma.

\begin{lemma}
  The space $Y_{m,n}(\sigma)$ is empty if and only if there exists an
  $i \leq m$ such that
  \[ 1 + n + i \leq \sigma(i). \]
\end{lemma}
\begin{proof}
  % Både dimensions-argumenter og stabiliserings-argumenter
  The exact condition is slightly strangely worded, but if there is
  such an $i$, then the first and last columns of the matrix
  \[ (e_{\sigma(i)},\dots,e_m,a_1,\dots,a_n,\sigma_1,\dots,\sigma_i) \]
  are linearly dependent, since
  \[ \sigma_i = e_{\sigma(i)} \]
  by definition. By counting the three different groups of vectors, we
  see that the matrix has
  \[ m - \sigma(i) + 1 + n + i \]
  columns. The condition on $i$ shows that
  \[ m - \sigma(i) + 1 + n + i \leq m - \sigma(i) + \sigma(i) = m, \]
  so the conditions for being in $Y_{m,n}(\sigma)$ can never be
  satisfied, no matter the choice of vectors $a_1,\dots,a_n$. This
  shows one direction.

  We will give two different proofs of the other direction. The first
  will proceed by a purely dimensional argument. Consider the two
  flags of interest,
  \begin{align*}
    \mathrm{Fl_R}(\Id) &= (V_1 \subset V_2 \subset \dots \subset V_m),
    \\
    \mathrm{Fl_L}(\sigma) &= (W_1 \subset W_2 \subset \dots \subset
    W_m).
  \end{align*}
  Assume that the condition on $\sigma$ and $n$ is satisfied, so for
  all $i \leq m$:
  \[ 1 + n + i > \sigma(i) \]
  In the flag formulation, this means that the subspaces $V_{m-n-k}$
  and $W_k$ intersect trivially for all $k$.

  We must now find the vectors $a_1,\dots,a_n$ and will do so from
  left to right. By looking at the matrix
  \[ \left( e_1, \dots, e_m, a_1, \dots, a_n, \sigma_1, \dots,
    \sigma_m \right), \]
  we see that we must choose $a_1$ so it is not in any of the
  following subspaces of $\C^m$:
  \begin{align*}
    a_1 &\not\in V_{m-1} \\
    a_1 &\not\in V_{m-n-1} \oplus W_1 \\
    \vdots & \\
    a_1 &\not\in V_1\oplus W_{m-n-1} \\
    a_1 &\not\in W_{m-n}
  \end{align*}
  But that means choosing $a_1$ in $\C^m$ after removing a
  finite number of proper subspaces, since the dimensions of the
  spaces on the right are $m-1$ for $V_{m-1}$ and $m-n$ for
  all the rest. This space is non-empty since $n\geq 1$, so this can
  always be done. We can continue this process inductively and get an
  element $(a_1,\dots,a_n)$ of $Y_{m,n}(\sigma)$, showing that the
  space is non-empty and concluding the proof.

  For good measure, we will also give an explicit element of
  $Y_{m,n}(\sigma)$.
  \fixme{Fortsæt her}
\end{proof}
\begin{remark}
  Note that if $n \geq m-1$, there is no such $i$ and the space
  $Y_{m,n}(\sigma)$ is non-empty. So the spaces become non-empty as
  soon as there are ``enough'' columns.
\end{remark}


\fixme{overvej ny sektion til gruppevirkning/kvotientrum}
The spaces have quite a few symmetries that could be of interest.
Since we are interested in linear independence of vectors in $\C^m$,
the group of complex units, $\T=(\C)^*$, acts on the spaces $Y_{m,n}$
in several ways. The 
most important for our purposes is the following:
\begin{definition}
  For any $i \in \set{1,\dots,n}$, the group $\T$ acts
  \fixme{husk at skrive $\T = \C^*$ i en intro et sted}
  on $Y_{m,n}(\sigma)$ by scaling the $i$'th column vector,
  using the formula
  \[ \l \cdot (a_1,\dots,a_{i-1},a_i,a_{i+1},\dots,a_n) =
  (a_1,\dots,a_{i-1},\l a_i,a_{i+1},\dots,a_n). \]
  Since the determinant of a matrix is linear in each column, the
  linear independence of the vectors is preserved.

  All of these actions commute with each other, so we get an action of
  $\T^n$ on $Y_{m,n}(\sigma)$:
  \[ (\l_1,\dots,\l_n) \cdot (a_1,\dots,a_n) = (\l_1 a_1,\dots, \l_n
  a_n). \]
\end{definition}

While all of these symmetries on an open subset of $\C^{mn}$ are nice
to work with, the spaces as defined above have some issues. In
particular, as we shall see in Theorem \ref{thm:forste}, the
fundamental group and the cohomology ring of
$Y_{m,n}(\sigma)$ both become larger as $n$ grows. This complicates
some of our arguments, but luckily it can be helped by taking the
quotient with the above group action. This gives us the spaces we are
actually interested in.

\begin{definition}
  For natural numbers $n$ and $m$ and invertible matrices $X$ and $Y$,
  we define the space $X_{m,n}(X,Y)$ as the quotient of $Y_{m,n}(X,Y)$
  using the group action above,
  \[ X_{m,n}(X,Y) = Y_{m,n}(X,Y) / \T^n. \]
\end{definition}

Note that since matrix multiplication is linear, it commutes with the
action of $\T^n$ on $Y_{m,n}(X,Y)$ and we get directly from Lemma
\ref{lem:rum-perm} that for each choice of $X,Y\in \GL_m(\C)$ there is
a permutation $\sigma \in S_m$ such that
\[ X_{m,n}(X,Y) \cong X_{m,n}(\sigma). \]
Hence most of our attention will be focused on these particular
spaces. The quotient spaces are naturally closely related to the
spaces $Y_{m,n}(X,Y)$. In particular, we have the following:

\begin{lemma}
  \label{lem:reduktion}
  There is a homeomorphism
  \[ Y_{m,n}(X,Y) \cong X_{m,n}(X,Y) \times \T^n. \]
\end{lemma}
\begin{proof}
  Define a map
  \[ X_{m,n}(X,Y) \to Y_{m,n}(X,Y) \]
  by mapping $A\cdot \T^n$ to the unique representative
  $\widetilde{A} \in Y_{m,n}(X,Y)$ with the last
  $n$ determinants all equal to one,
  \[ \Det_{X,Y}(\widetilde{A}) \in (\C^*)^{m-1} \times
  \{1\}^n. \]
  This map is continuous, as the composition with the quotient map,
  \[ \xymatrix{Y_{m,n}(X,Y) \ar[r]^q & X_{m,n}(X,Y) \ar[r] &
    Y_{m,n}(X,Y)}, \]
  is the map that uses the group action of $\T^n$ to change the last
  $n$ coordinates of $\Det_{X,Y}(A)$ to 1 and this is a continuous
  map. More precisely, it is given by the formula
  \[ (a_1,\dots,a_n) \mapsto \left(a_1,\dots,a_{m-1},\frac{1}{d_m}a_m,
    \frac{d_m}{d_{m+1}} a_{m+1},\dots,\hat{d}a_n\right), \]
  where $\Det_{X,Y}(a_1,\dots,a_n) = (d_1,\dots,d_{m+n-1})$ and the
  number $\hat{d}$ is an expression consisting of products and
  quotients of the last $n$ of these
  determinants, the exact terms of which depends on $n$ and $m$. If
  $\Det_{X,Y}(A) = (d_1,\dots,d_{m+n-1})$ denotes the determinants as
  above, the homeomorphism is the map
  \begin{align*}
    Y_{m,n}(X,Y) &\to X_{m,n}(X,Y) \times \T^n\\
    A &\mapsto (A\cdot \T^n, (d_m,\dots,d_{m+n-1})),
  \end{align*}
  with inverse
  \[ (A\cdot \T^n, (d_m,\dots, d_{m+n-1})) \mapsto
  \left(\overbrace{1,\dots,1}^{m-1},d_m,\frac{d_{m+1}}{d_m},\dots,
    \frac{1}{\hat{d}}\right) \cdot f(A). \]
\end{proof}

The above lemma is useful, since it oftens allows us to work with the
unquotiented spaces and still gain information about the
quotients. Note that in the definition of the homeomorphism
\[ Y_{m,n}(X,Y) \cong X_{m,n}(X,Y) \times \T^n , \]
 we might as well have
used the first $n$ determinants instead of the last, and we will
occasionally do this if it is easier for the arguments. \fixme{Tjek
  lige om vi faktisk gør dette, ellers så fjern det}
This lemma also immediately proves the claim that the fundamental
group and cohomology ring of $Y_{m,n}(X,Y)$ grows with $n$, by the
formulas
\begin{align*}
  \pi_1(Y_{m,n}(X,Y)) &\cong \pi_1(X_{m,n}(X,Y)) \times \Z^n, \\
  H^*(Y_{m,n}(X,Y)) &\cong H^*(X_{m,n}(X,Y)) \otimes H^*(\T^n).
\end{align*}
But more importantly, the last formula allows us to calculate the
cohomology of $X_{m,n}(X,Y)$ from the cohomology of $Y_{m,n}(X,Y)$, by
quotienting out the free groups coming from the cohomology ring of the
torus $H^*(\T^n)$. For example, we will later compute 
\[ H^1(X_{m,n}) \cong \Z^{m-1} \]
by instead computing
\[ H^1(X_{m,n}) \oplus \Z^n \cong H^1(Y_{m,n}) \cong \Z^{m+n-1}. \]

% Stabilisering
%% Flere retninger? Vi bruger ikke rigtig X_{m,n} -> X_{m+1,n}
%% Dette kunne måske linkes til Bruhat (så en introduktion til dette?)
%% Og definer grænsen, X_{m,\infty}
\section{Stabilisation} 
\label{sec:rum-stabil}

We would like to relate the spaces $Y_{m,n}(\sigma)$ that we obtain
for different values of
$m$ and $n$. Since an invertible, lower triangular matix $L$ preserves
the flag $\mathrm{Fl}_R(\Id)$, we will try to use this to define a map
\begin{align*}
  s(a_1,\dots,a_{n}) &= (La_1,\dots,La_{n},Le_1).
\end{align*}
At first, consider the spaces $Y_{m,n}$ and $Y_{m,n+1}$.
To ensure the right hand side actually is an element of $Y_{m,n+1}$,
we will start by considering the inverse of $s$, mapping a subset of
$Y_{m,n+1}$ to $Y_{m,n}$.
Consider an element $A=(a_1,\dots,a_{n+1})$ in $Y_{m,n+1}$ where
the entries of $a_{n+1}$ are all equal to one:
\fixme{Lav om på dette så det ser lidt bedre ud...}
\[ A = 
  \begin{pmatrix}
    \vrule & \vrule & & \vrule & \vspace{-0.5em}
                                 1 \\
    a_1 & a_2 & \dots & a_{n} & \vdots \\
    \vrule & \vrule & & \vrule & 1
  \end{pmatrix}. \] 
Then we can multiply through with the lower triangular, invertible
matrix $L^{-1}$ with ones on the diagonal and negative ones just below,
\fixme{Formuler bedre?}
\[ L^{-1} =
\begin{pmatrix}
   1 &  0 & 0 & \dots & 0 & 0 \\
  -1 &  1 & 0 & \dots & 0 &0 \\
   0 & -1 & 1 & \dots & 0 & 0 \\
   \vdots &  &  & \ddots & & \vdots \\
   0 & 0 & & & 1 & 0 \\
   0 & 0 & & & -1 & 1
\end{pmatrix}. \]\fixme{Få det til at se godt ud}
This reduces $A$ to
\[ L^{-1}A =
\begin{pmatrix}
  \vrule & \vrule & & \vrule & 
  1 \\
  \vrule & \vrule & & \vrule & \vspace{-0.5em}
  0 \\
  L^{-1}a_1 & L^{-1}a_2 & \dots & L^{-1}a_{n} & \vdots \\
  \vrule & \vrule & & \vrule & 0 \\
  \vrule & \vrule & & \vrule & 0
\end{pmatrix}, \]
and going a bit further we can see that
$(L^{-1}a_1,\dots,L^{-1}a_{n}) \in Y_{m,n}$ by computing
\[ L^{-1}\cdot \left(\Id_m,a_1,\dots,a_{n+1},\Id_m\right) =
\left(L^{-1},L^{-1}a_1,\dots,L^{-1}a_{n},
  \begin{array}{cccccc}
    1 & 1 & 0 & \dots & \dots & 0 \\
    0 & -1 & 1 & \dots & \dots& 0\\
    0 & 0 & -1 & \ddots & & \vdots\\
    \vdots & \vdots & & \ddots & 1 & 0 \\
    0 & 0 & \dots & & -1& 1
  \end{array}
\right). \] \fixme{Få til at se bedre ud og fylde mindre. hbox error}
This shows that $L^{-1}\cdot(a_1,\dots,a_{n})$ is in
$Y_{m,n}(L^{-1},U)$, where $U$ is the
upper triangular matrix given above. But this space is the same as
$Y_{m,n}$ since the right flag of a lower triangular matrix and the
left flag of an upper triangular matrix
is the same as the corresponding flag of the identity
matrix. \fixme{Henvisning hvis det er vist før}

Since the above is done by multiplying with a matrix, it commutes with
the action of the torus $\T^n$ on $Y_{m,n}$, and everything can be
repeated with $Y_{m,n}$ and $Y_{m,n+1}$ replaced by the spaces
$X_{m,n}$ and $X_{m,n+1}$.

The above considerations allow us define a map that relates the spaces
$X_{m,n}$ and $X_{m,n+1}$, by working with the inverse matrix $L$.
\begin{definition}
  \label{def:stabilisering}
  The \textit{stabilisation map} $s : X_{m,n} \to X_{m,n+1}$ is given
  by
  \[ s\left(a_1,\dots,a_{n-1}\right) =
  \left(La_1,\dots,La_{n-1},Le_1\right), \]
  where $L$ is the lower triangular matrix with every entry below the
  diagonal equal to 1,
  \[ L =
  \begin{pmatrix}
    1 & 0 & 0 & \dots & 0 \\
    1 & 1 & 0 & \dots & 0 \\
    1 & 1 & 1 & \dots & 0 \\
    \vdots & \vdots & \vdots & \ddots & \vdots \\
    1 & 1 & 1 & \dots & 1
  \end{pmatrix}. \]
\end{definition}

In the above, we only considered the spaces $X_{m,n}$ and $X_{m,n+1}$
defined by the identity permutation $\Id\in S_m$,
but there is a similar stabilisation for the other permutations. We
will again start by considering the unreduced spaces
$Y_{m,n}(\sigma)$. For
any permutation $\sigma$, the space $Y_{m,n+1}(\sigma)$ a union of
spaces homeomorphic to a torus times $Y_{m,n}(\tau)$, for various
permutations $\tau\in S_m$. The exact relation is given by the
following lemma.


\begin{lemma}
  \label{lem:permutation}
  Any choice of indices $I= (i_1<\dots<i_k) \subset \set{1,\dots,m}$
  defines a subspace of $Y_{m,n+1}(\sigma)$:
  \[ Y_{m,n+1}^I(\sigma) = \set{(a_1,\dots,a_{n+1}) \in
    Y_{m,n+1}(\sigma) \sdel 
    \begin{matrix} 
      (a_{n+1})_{i_j} \neq 0\; \forall i_j \in I, \\
      (a_{n+1})_j = 0 \;\forall j \not\in I
    \end{matrix} }. \]
  Here $(a_{n+1})_j$ denotes the $j$'th entry of the vector $a_{n+1}$.
  This space is empty if $I$ does not contain the number $\sigma(m)$,
  so we will assume that $\sigma(m) \in I$ but otherwise leave it out
  of the notation.
  
  These subspaces are disjoint and cover $Y_{m,n+1}(\sigma)$. There is
  a map $\varphi$,
  defined for an indexing set $I$ and a permutation $\sigma$, which
  gives a new permutation such that
  \[ Y_{m,n+1}^I(\sigma) \cong Y_{m,n+1}(\varphi(I,\sigma)) \times
  (\C^*)^{\size{I}}. \]
  If we define the permutation $\widehat \sigma = \sigma \cdot
  \big(m\; m-1\; \dots \; 1\big)$ and the indexing set is ${I = (i_1 <
  \dots < i_k)}$, then $\varphi$ is given by
  \[ \varphi\big(I,\sigma\big) = 
  \begin{cases}
    \begin{matrix} \widehat\sigma \end{matrix} & 
    \begin{array}{l} k = 0 \end{array} \\
    \begin{matrix}
      \big( i_k\; \sigma(m) \big) \cdot \varphi\left( (i_1 < \dots <
        i_{k-1}),\sigma\right) \\
    \end{matrix} &
    \begin{array}{l} i_k < \sigma(m) \text{ and } \\
      \widehat\sigma^{-1}(i_k) > \widehat\sigma^{-1}(i_r) \,
      \forall r < k\end{array}\\
    \begin{matrix}
      \varphi\left( (i_1 < \dots < i_{k-1}),\sigma\right)
    \end{matrix} &
    \begin{array}{l} \text{otherwise} \end{array}
  \end{cases}. \]
  Note that the index $\sigma(m)$ has once again been removed from the
  notation, 
  so the $k = 0$ case is really $I = (\sigma(m))$ and $i_k$
  is the last index not equal to $\sigma(m)$.
\end{lemma} \fixme{måske skal beviset skrives om. Det er muligvis
  rodet/for langt}
\begin{proof}
  Fix the permutation $\sigma \in S_m$. The observation that
  $Y^I_{m,n+1}(\sigma)$ is empty 
  if $\sigma(m)$ is not in $I$ follows from considering the matrix
  \[ (a_{n+1},\sigma_1,\dots,\sigma_{m-1}) =
  (a_{n+1},e_{\sigma(1)},\dots,e_{\sigma(m-1)}). \]
  For this matrix to be invertible we must have $(a_{n+1})_{\sigma(m)}
  \neq 0$. If $\sigma(m)$ is not in $I$, then the conditions on
  $Y^{I}_{m,n+1}(\sigma)$ can never be satisfied and the space must be
  empty. We will from now on assume that $I$ contains $\sigma(m)$.
  
  To see that the spaces are disjoint and cover all of
  $Y_{m,n+1}(\sigma)$, we can see that an element
  $(a_1,\dots,a_{n+1})\in Y_{m,n+1}(\sigma)$ is in exactly one
  subspace, namely the one corresponding to $I = \set{ j \in
    \set{1,\dots,m} \mid (a_{n+1})_j \neq 0}$.
  
  To define the function $\varphi$ and the homeomorphism
  \[ Y^I_{m,n+1}(\sigma) \cong Y_{m,n}(\varphi(I,\sigma))\times
  (\C^*)^{\size{I}}, \]
  we will proceed as in the previous case. The map to the torus
  $(\C^*)^{\size{I}}$ is given by remembering the non-zero entries of
  $a_{n+1}$, so we can immediately reduce to the case where all
  entries of $a_{n+1}$ are either zero or one.

  An element of $Y_{m,n+1}^I(\sigma)$ consists of vectors
  $(a_1,\dots,a_{n+1})$ with the matrix
  \[ M = (a_{n+1},\sigma_1,\dots,\sigma_{m-1}) \]
  invertible. It is this matrix that we want to factor as
  $L\cdot \varphi(I,\sigma)\cdot U$, where $L$ is an invertible, lower
  triangular matrix and $U$ is an invertible, upper triangular
  matrix. To do this, we want to make row operations on the matrix
  corresponding to multiplying with $L^{-1}$ until we end up with an
  upper triangular matrix with permuted rows, corresponding to
  $\varphi(I,\sigma) \cdot U$. Multiplying with $L^{-1}$ corresponds
  to making row operations where each row is only allowed to affect
  the rows below it. We want to remove all non-zero entries in the
  first column except for one, and we want this single entry to have
  the lowest possible index. We also want to make the matrix as close
  as possible to upper triangular, which is done by making sure that
  each operation does not introduce a non-zero entry to the left of an
  already existing non-zero entry.
  
  If $k = 0$, an element $Y^I_{m,n+1}(\sigma)$ looks like
  \[ (a_1,\dots,a_n,\sigma_m). \]
  Since we are in $Y_{m,n+1}(\sigma)$, the linear independence
  condition gives us that
  \[ (a_1,\dots,a_n) \in Y_{m,n}(\tau), \]
  where $\tau$ is the permutation
  \[ \tau(k) =
  \begin{cases}
    \sigma(k-1) & k > 1 \\
    \sigma(m) & k = 1
  \end{cases}. \]
  We recognize this as the permutation $\widehat\sigma$ and get the
  formula
  \[ \varphi((\sigma(m)),\sigma) = \widehat\sigma. \]
  The homeomorphism is in this case given by
  \begin{align*}
    Y_{m,n+1}^{(\sigma(m))}(\sigma) &\to Y_{m,n}(\widehat\sigma)\times
                                      \C^* \\
    (a_1,\dots,a_n,\lambda\cdot\sigma_m) &\mapsto
                                           (a_1,\dots,a_n,\lambda)
  \end{align*}
  
  Now consider the case $k > 0$, so $I = (i_1 < \dots < i_k)$. If $i_k
  > \sigma(m)$, we can write
  $(a_{n+1},\sigma_1,\dots,\sigma_{m-1})$ in block form as
  \[
  \begin{pmatrix}
    I_1 & A \\
    1 & 0 \\
    I_2 & B \\
    1 & C \\
    0 & D
  \end{pmatrix}.
  \]
  $I_1$ is a vector with ones on the entries corresponding
  to $\set{i_j \mid i_j < \sigma(m)}$ and zeroes elsewhere, while
  $I_2$ is similar except we consider indices greater than
  $\sigma(m)$. The first 1 is on the $\sigma(m)$'th row and the next
  is on the $i_k$'th row. All entries below are zero. $A$, $B$, $C$
  and $D$ are the matrices consisting of the corresponding rows from
  $(\sigma_1,\dots,\sigma_{m-1})$. But by multiplying with a lower
  triangular matrix, we can reduce this matrix to
  \[
  \begin{pmatrix}
    I_1 & A \\
    1 & 0 \\
    I_2 & B \\
    0 & C \\
    0 & D
  \end{pmatrix}.
  \]
  But this shows that $\varphi(I,\sigma) =
  \varphi\big((i_1<\dots<i_{k-1}),\sigma\big)$ when $i_k >
  \sigma(m)$.

  By possibly using the above multiple times, we reduce to the case
  where $i_k < \sigma(m)$. Now we would like to reduce the length of
  $I$ by one, and consider how this affects the resulting permutation.
  Consider $\widehat\sigma^{-1}(i_k)$. If
  there is a $j < k$ with $\widehat\sigma^{-1}(i_k) <
  \widehat\sigma^{-1}(i_j)$, then we can write
  $(a_{n+1},\sigma_1,\dots,\sigma_{m-1})$ as
  \[
  \begin{pmatrix}
    \vdots & \\
    1 & \dots & \dots & \dots & 1 & \dots\\
    \vdots \\
    1 & \dots & 1 & \dots & \dots & \dots\\
    \vdots & \\
    1 &  \dots & \dots & \dots & \dots & \dots\\
    \vdots
  \end{pmatrix}. \]
  In the above, only the $i_j$'th, $i_k$'th and $\sigma(m)$'th rows
  have been written out, and only the non-zero entries. But this
  matrix can be reduced to
  \[
  \begin{pmatrix}
    \vdots & \\
    1 & \dots & \dots & \dots & 1 & \dots\\
    \vdots \\
    0 & \dots & 1 & \dots & -1 & \dots\\
    \vdots & \\
    1 & \dots & \dots  & \dots & \dots & \dots\\
    \vdots
  \end{pmatrix}. \]
  We are now unable to change row $i_k$, since there are no non-zero
  entries above the left-most one. The additional $-1$ that was
  introduced can be removed by multiplying on the right with an upper
  triangular matrix that subtracts column $\widehat\sigma^{-1}(i_k)$
  from column $\widehat\sigma^{-1}(i_j)$. So we are able to remove the
  one in the $i_k$'th row without changing the resulting permutation,
  showing that $\varphi\big(I,\sigma\big) =
  \varphi\big((i_1<\dots<i_{k-1}),\sigma\big)$ in this case.

  Now assume that there is no such $j$. Then the row $i_k$ is our best
  candidate for removing the bottom-most one in row $\sigma(m)$. If we
  do this row operation, by multiplying with $\widehat L$, we get
  \[ \varphi(I,\sigma)\cdot U = L^{-1}\widehat{L} \cdot M \]
  But when we perform row operations on $\widehat{L} \cdot M$, we are
  only changing rows with index less than or equal to $i_{k-1}$, since
  this is the last non-zero entry in the first column. Hence $L^{-1}$
  has the form
  \[ L^{-1} = 
  \begin{pmatrix}
    \vdots \\
    * & \dots & * & 1 & \dots & \dots & \dots \\
    \vdots \\
    0 & \dots & \dots & \dots & \dots & 1 & \dots \\
    \vdots
  \end{pmatrix},
% \left(
%   \begin{array}{c|c}
%     L' & 0 \\
%     \hline
%     0 & \Id
%  \end{array} \right)
  \]
  where we once again write out row $i_k$ and $\sigma(m)$, and use the
  stars to indicate entries that may be non-zero but that are not
  important.
  If we switch row $i_k$ and row $\sigma(m)$ along with column
  $i_k$ and column $\sigma(m)$, we get another lower triangular
  matrix:\fixme{beslut (overalt) om der bruges , eller ; i
    permutationer}
  \[ (i_k \; \sigma(m)) L^{-1} (i_k \; \sigma(m)) = 
  \begin{pmatrix}
    \vdots \\
    0 & \dots & \dots & 1 & \dots & \dots & \dots \\
    \vdots \\
    * & \dots & * & \dots & \dots & 1 & \dots \\
    \vdots
  \end{pmatrix}.\]
  Inserting in our original equation, we see
  \[ (i_k \; \sigma(m)) \varphi(I,\sigma) \cdot U = L^{-1} \cdot (i_k \;
  \sigma(m)) \cdot \widehat{L} M. \]
  But $(i_k \; \sigma(m)) \cdot \widehat{L} M$ is just $M$ with the
  one in the first column, row $i_k$, replaced by zero. This gives the
  formula
  \[ (i_k \; \sigma(m)) \cdot \varphi(I,\sigma) =
  \varphi\big((i_1<\dots<i_{k-1}),\sigma\big) \]
  and our desired formula for $\varphi$ follows by multiplying over.
\end{proof}
\fixme{Sørg for at permutations-notationen er beskrevet,
  e.g. $\sigma_1 = e_{\sigma(1)}$.}


\fixme{Noget afsluttende}

\section{The limit space}
\label{sec:rum-gr}

Using the stabilisation defined previously, we get a directed system
of topological spaces,
\[ \xymatrix{X_{m,1} \ar[r]^s & X_{m,2} \ar[r]^s & \dots \ar[r]^s &
  X_{m,n} \ar[r]^s & X_{m,n+1} \ar[r]^s & \dots } \]
When considering such a system, it can sometimes be helpful to
take the limit as $n$ approaches infinity.

\begin{definition}
  The space $X_{m,\infty}$ is defined as the direct limit of the
  directed system of spaces,
  \[ X_{m,\infty} = \varinjlim_n X_{m,n}. \]
  As a set, $X_{m,\infty}$ is the disjoint union of all the spaces,
  where we identify $A \in X_{m,k}$ with $B \in X_{m,k'}$ if they
  eventually map to the same thing under $s$:
  \[ A \sim B \iff \exists\, i,j : s^i(A) = s^j(B). \]
  We give $X_{m,\infty}$ the finest topology such that all the
  maps 
  \[ \imath_n : X_{m,n} \to X_{m,\infty} \]
  are continuous. So a subset $U \subset X_{m,\infty}$ is open if and
  only if $\imath_n^{-1}(U)$ is open in $X_{m,n}$ for all n. \fixme{Tjek
    continuous alle steder...}
\end{definition}

We could just as well define the space \fixme{Overvej lige om det her
  er rigtigt og om det bare skal fjernes igen}
\[ Y_{m,\infty} = \varinjlim_n Y_{m,n} \]
and note that all stabilisation maps are linear. This means we get an
action on $Y_{m,\infty}$ of the infinite torus,
\[ \T^\infty = \varinjlim_n \T^n, \]
defined by taking the
limit of the
inclusion maps $\T^n \to \T^{n+1}$ that inserts a one on the last
coordinate.
Then we could define the limit space as
\[ X_{m,\infty} = Y_{m,\infty}/\T^\infty. \]

% Noget med at $s^m \simeq \Id$?


% Hvis ikke introduceret endnu ville spektralfølgen give mening
% her. Hvis allerede introduceret i det generelle tilfælde ville det
% nok give mening at specialisere til Y_{m,n} eller X_{m,n} her og
% beskrive de forskellige søjler.
%% Her ville det måske også give mening at se på H^1 og vise at den
%% giver det "rigtige" svar, \Z^{m-1}, for at vise at følgen kan
%% bruges til noget.

\fixme{Flyt dette hen hvor det giver mening. Muligvis et andet
  kapitel.}

\section{The first cohomology of $X_{m,n}$}
\label{sec:firstcohom}

We would like to use the spectral sequence to compute the first
cohomology group of $X_{m,n}$ for any $m$ and $n$. As discussed at the
end of Section \ref{sec:rum}, we will do this by computing
$H^1(Y_{m,n})$ and
then using Lemma \ref{lem:reduktion} to calculate the cohomology of
$X_{m,n}$ via factoring out the extra free group from the torus
in the isomorphism
\[ H^1(Y_{m,n}) \cong H^1(X_{m,n}) \oplus H^1(\T^n). \]

From the theory
underlying the spectral sequence, see e.g. \cite[Chapter
1]{hatcherss}, we get that there is a subgroup
\[ 0 \subset H \subset H^1(Y_{m,n}) \]
and isomorphisms
\begin{align*}
  E_\infty^{0,1} &\cong H^1(Y_{m,n})/H, \\
  E_\infty^{1,0} &\cong H.
\end{align*}
So to calculate $H^1(Y_{m,n})$, we have to calculate the two groups
$E_\infty^{1,0}$ and $E_\infty^{0,1}$ and see what groups fit
into the short exact sequence
\[ 0 \to E^{1,0}_\infty \to H^1(Y_{m,n}) \to E^{0,1}_\infty \to 0. \]

Drawing the relevant part of the first page of the spectral sequence,
we are considering the diagram illustrated in Figure
\ref{fig:foerste}.
\begin{figure}[ht]
  \[ \xy
  (5,4); (5,-13); **\dir{-}; (5,4); *\dir{>};
  (5,-13); (73,-13); **\dir{-}; *\dir{<};
  (0,-15)*{q}; (5,-17)*{p};
  (1,-17); (4,-14); **\dir{-};
  (0,0)\xymatrix@R=1em@C=1.7em{
    1 & H^1(F_0) \ar[r]& H^2(F_1,F_0) \ar[r] & H^3(F_2,F_1) \\
    0 & H^0(F_0) \ar[r] & H^1(F_1,F_0) \ar[r] & H^2(F_2,F_1) \\
    & 0 & 1 & 2
  } \endxy \]
  \caption{The relevant part of $E_1$ for computing $H^1(Y_{m,n})$.}
  \label{fig:foerste}
\end{figure}
But by writing out our specific case, we can see that the the groups
$H^1(F_1,F_0), H^2(F_2,F_1)$ and $H^3(F_2,F_1)$ are all zero, so the
diagram simplifies to Figure \ref{fig:anden}. \fixme{Gør dette
  tidligere}
\begin{figure}[ht]
  \[ \xy
  (5,4); (5,-13); **\dir{-}; (5,4); *\dir{>};
  (5,-13); (75,-13); **\dir{-}; *\dir{<};
  (0,-15)*{q}; (5,-17)*{p};
  (1,-17); (4,-14); **\dir{-};
  (0,0)\xymatrix@R=1em@C=1.7em{
    1 & H^1(Y_{m,n-1}\times \T^m) \ar[r]& H^2(F_1,F_0) & 0 \\
    0 & H^0(Y_{m,n-1}\times \T^{m}) & 0 & 0 \\
    & 0 & 1 & 2
  } \endxy \]
  \caption{The simplified $E_1$.}
  \label{fig:anden}
\end{figure}
From the figure, we can see that the group $E_\infty^{1,0}$
must be zero, so the exact sequence we are considering gives an
isomorphism
\[ \xymatrix{0 \ar[r] & H^1(Y_{m,n}) \ar[r]^>>>>>{\cong} & E^{0,1}_\infty
  \ar[r] & 0}. \]

By doing explicit calculations for small $m$ and $n$, we are led to
the following theorem. \fixme{Inkluder disse/henvis til eksemplerne i
  starten}

\begin{theorem}
  \label{thm:forste}
  The first cohomology group of $Y_{m,n}$ is isomorphic to $\Z^{m+n-1}$.
\end{theorem}
\begin{proof}
  We will proceed by induction on $n$. By looking at Example
  \ref{ex:n=1}, we see that
  \[ H^1(Y_{m,1}) \cong H^1(\T^{m}) \cong \Z^m = \Z^{m+1-1} \]
  as desired.

  Now assume the theorem is correct for 
  \[ H^1(Y_{m,n-1}) \cong \Z^{m+n-2}. \] \fixme{Bør vise at $Y_{m,n}$
    0-sh tidligere}
  By excision, we know that the group $H^2(F_1,F_0)$ is given by
  \[ H^2(F_1,F_0) \cong \bigoplus_{i=1}^{m-1} H^0(Y_{m,n-1}\big((i\;
  i+1)\big) \otimes H^2(D,D-0) \cong \Z^{m-1}. \]\fixme{Sørg for
    at vi ved dette}
  \begin{figure}[ht]
    \[ \xy
    (5,4); (5,-13); **\dir{-}; (5,4); *\dir{>};
    (5,-13); (61,-13); **\dir{-}; *\dir{<};
    (0,-15)*{q}; (5,-17)*{p};
    (1,-17); (4,-14); **\dir{-};
    (0,0)\xymatrix@R=1em@C=1.7em{
      1 & \Z^{m+n-2}\oplus\Z^m \ar[r]& \Z^{m-1} & 0 \\
      0 & \Z & 0 & 0 \\
      & 0 & 1 & 2
    } \endxy \]
    \caption{$E_1$, assuming the induction hypothesis.}
    \label{fig:tredje}
  \end{figure}
  This simplifies Figure \ref{fig:anden} to Figure
  \ref{fig:tredje}. Since the groups involved will never change after
  turning the page to $E_2$, we see that we only have to compute the
  group
  \[ E_\infty^{0,1} = E_2^{0,1} = \ker \left( d_1 : E_1^{0,1} \to
    E_1^{1,1} \right). \]
  The differential $d_1$ is given by the differential in the long
  exact sequence for the pair $(F_1,F_0)$. By studying the long exact
  sequence, the differential is given by the composition
  \fixme{Beskriv tidligere?}
  \[ d_1 : \xymatrix{ H^1\left(Y_{m,n-1}\times \T^m\right)
    \ar[r]^<<<<<<<<<{\oplus\imath^*}  & \bigoplus_{i=1}^{m-1}
    H^1\left(Y_{m,n-1}\big((i\; i+1)\big)\times \T^{m-1} \times
      (D-0)\right)  \\
    \hspace{5em}\ar[r]^<<<<<{\oplus\delta} & \bigoplus_{i=1}^{m-1}
    H^2\left(Y_{m,n}\big((i\; i+1)\big)\times \T^{m-1} \times
      (D,D-0)\right).}\] \fixme{Få pilen til at sidde ordentligt}
  The first map is induced by the inclusions
  \[ \imath : Y_{m,n-1}\big((i\; i+1)\big) \times \T^{m-1} \times (D-0)
  \to Y^{m,n-1}\times \T^m, \] 
  given by first using the inverse of the homeomorphism from Lemma
  \ref{lem:permutation}, given by the indices $I = (1 < 2 < \dots <
  i-1 < i+1 < \dots m)$,
  \[ Y_{m,n}^I \to Y_{m,n+1}\big((i \; i+1)\big) \times \T^{m-1}, \]
  and then inserting the coordinate from $(D-0)$ on the $i$'th
  coordinate of the last column. \fixme{Beskriv at $D-0$ er meget
    lille et sted tidligere}
  The second map consists of the differentials from the long exact
  sequences of the pairs
  \[ Y_{m,n-1}\big((i\; i+1)\big) \times \T^{m-1}\times (D,D-0). \]
  We want to see that $d_1$ is surjective, so we will compute the
  image of the element 
  \[ \alpha_j = 1\otimes\dots 1\otimes [T] \otimes 1 \otimes \dots
  \otimes 1 \in H^1(Y_{m,n-1}\times \T^{m}),\] 
  where the class $[T]\in H^1(T)$ is in the $j$'th tensor coordinate.
  We will consider each inclusion separately. If $i \neq j$, the
  inclusion will not touch the $j$'th coordinate of $\T^m$, so on
  cohomology classes it will map $\alpha_j$ to an element that is of
  the form
  \[ \imath^*(\alpha_j) = \widetilde{\alpha_j}\otimes 1 \in
  H^*(Y_{m,n-1}\big((i\; i+1) \big) \times \T^{m-1}) \otimes
  H^*(D-0) \]
  And since $\delta$ is given by \fixme{Bør nok også argumenteres for
    et sted}
  \[ \delta(\beta) =
  \begin{cases}
    \widetilde{\beta}\otimes [(D,D-0)] & \text{if } \beta =
    \widetilde{\beta}\otimes [D-0] \\
    0 & \text{otherwise}
  \end{cases}, \]
  we see that $\delta(\imath^*(\alpha_j)) = 0$ in this case. If $i =
  j$, then the same considerations give us that
  \[ \delta\imath^*(\alpha_j) = 1 \otimes [(D,D-0)] \in
  H^2(Y_{m,n-1}\big((i\; i+1) \big) \times \T^{m-1}\times (D,D-0)). \]
  The group $H^2(F_1,F_0)$ are generated by exactly these elements, so
  the map $d_1$ is surjective. Since the groups involved are all free,
  we know the kernel is also free and we can calculate the rank as the
  rank of the domain minus the rank of the codomain, giving us
  \[ H^1(Y_{m,n}) \cong E_{\infty}^{0,1} = \ker d_1 \cong
  \Z^{m+n-2+m-(m-1)} = \Z^{m+n-1}. \]
\end{proof}

An immediate consequence is the following result.

\begin{corollary}
  The first cohomology group of $X_{m,n}$ is $\Z^{m-1}$.
\end{corollary}
\begin{proof}
  By using the Künneth formula and the previous theorem, we get an
  isomorphism \fixme{Henvisning til kunneth?}
  \[ H^1(X_{m,n}) \oplus H^1(\T^n) \cong  H^1(Y_{m,n}) \cong
  \Z^{m+n-1}. \]
  Since any torsion in $H^1(X_{m,n})$ would show up in $H^1(Y_{m,n})$,
  it must be free. We can calculate the rank from the above and get
  \[ H^1(X_{m,n}) \cong \Z^{m-1}. \]
\end{proof}

To make this abstract isomorphism slightly more concrete, we will
realise it as a map between spaces. Consider the image of an element
$v \in Y_{m,1}$ in $Y_{m,n}$ under the stabilisation map. This has the
form 
\[ s^{n-1}(v) = s^{n-1}
\begin{pmatrix}
  v_1 \\
  \vdots \\
  v_m
\end{pmatrix} =
\begin{pmatrix}
  \vrule & \vrule & & \vrule \\
  L^{n} v & L^{n-1}e_1 & \dots & Le_1 \\
  \vrule & \vrule & & \vrule
\end{pmatrix},
\]
where $L$ is the lower triangular matrix that defines the
stabilisation map, see Definition \ref{def:stabilisering}. Computing the
determinants gives
\[ \Det(s^{n-1}(v)) = (\pm v_1,\dots,\pm v_m, \pm 1, \dots,\pm 1), \]
with the signs dependent on $n$ and $m$. Using the action of
$\T^{n}$ to change the last $n-1$ determinants as in
\fixme{find ud af hvad der skal refereres til
  her}, we can get an element of $Y_{m,n}$ with any prescribed
determinant. Altogether this defines a lift, $f$, of the
determinant map,
\[ \xymatrix{ \T^{m+n-1} \ar[r]^f \ar@/_1.5pc/[rr]_\Id & Y_{m,n}
  \ar[r]^<<<<<\Det & \T^{m+n-1}}. \]
By taking first cohomology, we get maps
\[ \xymatrix{ H^1(Y_{m,n}) \ar@/^/[r]^{f^*} & H^1(\T^{m+n-1})
  \ar@/^/[l]^{\Det^*} }, \]
with 
\[ f^* \circ \Det^* = \Id. \]
This shows that $\Det^*$ is injective and $f^*$ is surjective. But by
the previous theorem, $f^*$ is a map between free abelian groups of
the same rank, hence it must also be injective and so an
isomorphism with inverse $\Det^*$. This shows that the determinant map
is an isomorphism on the first cohomology groups.
This result immediately extends to the quotiented spaces, with
\fixme{Sørg for at $\TT{m+n-1}{n}$ er introduceret før. Det er ikke
  den oplagte virkning...}
\[ \Det^* : H^1(\TT{m+n-1}{n}) \to H^1(X_{m,n}) \]
an isomorphism. Furthermore it shows that the stabilisation maps on
degree one cohomology are surjective, since the stabilisation maps
\[ s: Y_{m,n} \to Y_{m,n+1} \]
change the signs of some determinants and add an extra determinant
that is always plus or minus one. So if the class $d_i \in
H^1(Y_{m,n+1})$ represents the image of moving once around zero in the
$i$'th coordinate of $\T^{n+1}$, then
\[ s^*(d_i) =
\begin{cases}
  \pm d_i & 1 \leq i \leq n \\
  0 & \text{otherwise}
\end{cases} \] \fixme{Muligvis skift så vi bruger $x_i$ eller $t_i$ eller
  whatever der nu giver mening når jeg har skrevet kapitlet med
  udregninger}


% The theorem and the corollary illustrate the problems with using
% $Y_{m,n}$ instead of $X_{m,n}$, even though the space $Y_{m,n}$ is
% conceptually simpler. We want to consider the limit spaces
% $X_{m,\infty}$ and $Y_{m,\infty}$, but if we compute the first
% homology group of the direct limit using e.g. \cite[\S14.6]{may}, we
% see
% \begin{align*}
%   H_1(Y_{m,\infty}) &= H_1(\varinjlim_n Y_{m,n}) = \varinjlim_n
%   H_1(Y_{m,n}) = \varinjlim_n \Z^{m+n-1} = \Z^\infty \\
%   H_1(X_{m,\infty}) &= H_1(\varinjlim_n X_{m,n}) = \varinjlim_n
%   H_1(X_{m,n}) = \varinjlim_n \Z^{m-1} = \Z^{m-1}.
% \end{align*}
% The reason we compute homology instead of cohomology is that
% homology commutes with direct limits. . So for $Y_{m,\infty}$ the first homology group is not finitely generated.

%%%%% Previous
% \chapter{The space $\X$}

% % Snak snak. Måske deles op i to?
% % Noget om starten, hvorfor dette er interessant. Se bagside af blok
% % Også noget om rummene, definition, symmetrier (specielt X_{m,n}
% % \cong X_{n,m}, måske filtreringen?
% This chapter will introduce the spaces to be studied, along with a
% brief listing of various useful structure.

% \begin{definition}
%   For natural numbers $m$ and $n$, the space $X_{m,n} \subset
%   \C^{mn}$ is defined as
%   \[ X_{m,n} = \set{(a_1,\dots,a_n) \in (\C^m)^n \delim\Bigg\vert\delim
%     \begin{matrix}
%       \text{Any } m \text{ subsequent vectors in } \\
%       (e_1,\dots,e_m,a_1,\dots,a_n,e_1,\dots,e_m) \\
%       \text{ are linearly independent.}
%     \end{matrix} } \]
%   For any two elements $X = [x_1,\dots,x_m]$ and $Y = [y_1,\dots,y_m]$
%   in $\GL_m$, define the space $X_{m,n}(X,Y)$ as
%   \[ X_{m,n}(X,Y) = \set{(a_1,\dots,a_n) \in (\C^m)^n
%     \delim\Bigg\vert\delim
%     \begin{matrix}
%       \text{Any } m \text{ subsequent vectors in } \\
%       (x_1,\dots,x_m,a_1,\dots,a_n,y_1,\dots,y_m) \\
%       \text{ are linearly independent.}
%     \end{matrix} } \]
%   The special case $X=\Id$ will be denoted $\X(Y)$.
%   Elements $(a_1,\dots,a_n)\in \X$ will be identified
%   with the $m\times n$ matrix $A = [a_1,\dots,a_n]$ without mention.
% \end{definition}

% Since the space $\X(X,Y)$ only depends on linear independence of
% subsequent vectors, it can be described by giving the two flags in
% $\C^m$ defined from the columns of the matrices
% $X=\left[x_1,\dots,x_m\right]$ and $Y=\left[y_1,\dots,y_m\right]$:
% \begin{align*}
%   \mathrm{Fl_R}(X) &= \Big(\spa(x_m) \subset \spa(x_{m-1},x_m) \subset
%     \dots \subset \C^m\Big) \\
%   \mathrm{Fl_L}(Y) &= \Big(\spa(y_1) \subset \spa(y_1,y_2) \subset
%     \dots \subset \C^m\Big)
% \end{align*}
% Multiplying with $X^{-1}$ on each vector in $\X(X,Y)$ defines a
% homeomorphism 
% \[ \X(X,Y)\cong \X(\Id,X^{-1}Y) = \X(X^{-1}Y) \]
% Hence we only need to consider these spaces.

% Note that we can also define the space $X_{m,n}$ as the preimage of
% $(\C^*)^{m+n-1}$ under the map $\Det : (\C^m)^n \to \C^{m+n-1}$,
% given by
% \begin{align*}
%   \Det(a_1,\dots,a_n) = \Big( &\det(e_2,\dots,e_m,a_1),\\
%   &\det(e_3,\dots,e_m,a_1,a_2),\\
%   &\dots, \\
%   &\det(a_n,e_1,\dots,e_{m-1}) \Big) 
% \end{align*}
% where $\det$ is the determinant map. There is a similar description
% of $X_{m,n}(Y)$ as the preimage under a map $\Det_{Y}$. Since the map
% $\Det_Y$ is continuous, this shows that the space $\X(Y)$ is open in
% $\C^{mn}$.

% \section{Stabilization}

% We would like to relate the spaces $X_{m,n}$ and $X_{m,n+1}$. To do
% this, consider a lower-triangular $m\times m$ invertible matrix
% $L$. This preserves the right flag $\mathrm{Fl_R}(\Id)$, so it defines
% a homeomorphism $\tilde L : \X(Y) \to \X(LY)$ by
% \begin{align*}
%   \tilde L(a_1,\dots,a_n) &= (La_1,\dots,La_n)
% \end{align*}

% By using the Bruhat decompositon of the general linear group,
% any matrix $Y$ can be written uniquely as a product $L \sigma U$,
% where $L$ is an invertible lower triangular matrix, $U$ is an
% invertible upper triangular matrix, and $\sigma$ is a
% permutation matrix, see \cite[Example~1.2.11]{bjorner} or
% \cite[Proposition~4.5]{hiller}. Since the flag $\mathrm{Fl_L}(\sigma
% U)$ is preserved by right multiplication with upper triangular
% matrices, this only depends on the permutation $\sigma$ and not on
% $U$. All of this together shows the following lemma:

% \begin{lemma}
%   For $Y\in \GL_m$ with $Y = L\sigma U$, the space $\X(Y)$ is
%   homeomorphic to $\X(\sigma)$.
% \end{lemma}

% Using this lemma, we only need to understand the spaces
% $\X(\sigma)$ to understand all of the spaces $\X(Y)$. More importantly,
% it turns out that the space $X_{m,n+1}(\sigma)$ is given as a union of
% spaces homeomorphic to a torus times $\X(\tau)$, for various
% permutations $\tau$. This is summarized in the following lemma.

% \begin{lemma}
%   Any choice of indices $I= (i_1<\dots<i_k) \subset \set{1,\dots,m}$
%   gives a subspace of $X_{m,n+1}(\sigma)$:
%   \[ X_{m,n+1}^I(\sigma) = \set{(a_1,\dots,a_{n+1}) \in
%     X_{m,n+1}(\sigma) \sdel 
%     \begin{matrix} 
%       (a_{n+1})_{i_j} \neq 0\; \forall i_j \in I, \\
%       (a_{n+1})_j = 0 \;\forall j \not\in I
%     \end{matrix} } \]
%   This space is empty if $I$ does not contain the number $\sigma(m)$,
%   so we will assume that $\sigma(m) \in I$ but otherwise leave it out
%   of the notation.
  
%   These subspaces cover $X_{m,n+1}(\sigma)$. There is a map $\varphi$,
%   defined for an indexing set $I$ and a permutation $\sigma$, which
%   gives a new permutation such that
%   \[ X_{m,n+1}^I(\sigma) \cong \X(\varphi(I,\sigma)) \times
%   (\C^*)^{\size{I}} \]
%   If we define the permutation $\widehat \sigma = \sigma \cdot
%   \big(m\; m-1\; \dots \; 1\big)$ and the indexing set is $I = (i_1 <
%   \dots < i_k)$, then $\varphi$ is given by
%   \[ \varphi\big(I,\sigma\big) = \left(\prod_{j \in J} \big(
%       i_j \; \sigma(m) \big) \right) \widehat \sigma \]
%   where the product is over the set
%   \[ J = \set{ j \in\set{1,\dots,k} \delim\mid\delim i_j <
%     \sigma(m),\; \widehat\sigma^{-1}(i_j) > \widehat\sigma^{-1}(i_r) \,
%     \forall r < j } \]
% \end{lemma}

% The proof of this lemma relies on taking $X_{m,n+1}^I(\sigma)$ and
% multiplying with a lower triangular matrix to reduce it to something
% in $\X(\tau)$. By considering the matrix used, the formula
% for $\tau = \varphi(I,\sigma)$ given in the lemma can be found. In
% practice, it is
% often easier to do the reduction by hand rather than using the lemma,
% but the formula can be useful. For example, it is clear that
% $\varphi(I,\sigma)$ is given as a product of transpositions
% and $\varphi\big((\sigma(m)), \sigma\big)$. By
% considering exactly which transpositions, we can show that for all $I$
% we have
% \[ \varphi(I,\sigma) \leq \varphi\big((\sigma(m)), \sigma\big) \] 
% in the Bruhat ordering on the symmetric group. The Bruhat
% order is defined in \cite{bjorner}. It also follows that for all
% $\sigma$, 
% \[ \varphi\big((1<2<\dots<m),\sigma\big) = \Id \]
% With this we can identify $X_{m,n}$ with the subspace
% of $X_{m,n+1}$ where the last column has a 1 on each entry:
% \[ X_{m,n} \cong \set{ (a_1,\dots,a_{n+1}) \in X_{m,n+1} \mid
%   a_{n+1} = (1,\dots,1) } \]
% The identification is given by the map $s : \X \to X_{m,n+1}$,
% \[ s(a_1,\dots,a_n) = (La_1,\dots,La_n,Le_1) \]
% where $L$ is the lower triangular matrix with all entries on or below
% the diagonal equal to one:
% \[ L =
% \begin{pmatrix}
%   1 & 0 & \dots & 0 \\
%   \vdots & \ddots & \ddots & \vdots \\
%   \vdots &  & \ddots & 0 \\
%   1 & \dots & \dots & 1
% \end{pmatrix} \]

% These maps give a directed system of spaces,
% \[ \xymatrix{ X_{m,1} \ar[r]^s & X_{m,2} \ar[r]^s & \dots \ar[r]^s &
%   X_{m,n} \ar[r]^s & X_{m,n+1} \ar[r]^s & \dots } \]
% and we can take the direct limit of this system,
% \[ X_{m,\infty} = \varinjlim_n \X \]
% The direct limit is the disjoint union of all the spaces $\X$, with the
% identification $A \sim B$ if there are $j,k$ such that
% $s^k(A) = s^j(B)$. The topology is the finest such that the maps $\X
% \to X_{m,\infty}$ are all continuous.

% \section{Symmetries}

% % Skriv noget her og byt rundt på rækkefølge?
% % Noget om virkninger af \C^* på \X, e.g. virk på søjler, rækker,...
% There are various symmetries and group actions worth keeping in mind
% when working with these spaces. Some of these will be described here.

% Since we are interested in linear independence of vectors in $\C^m$,
% the group of complex units $\C^*$ acts on $\X(\sigma)$ in various
% ways. For any index $i$, we can define a group action $l_i$ by scaling
% the $i$\textsuperscript{th} vector:
% \begin{align*}
%   l_i : \C^* \times \X(\sigma) &\to \X(\sigma) \\
%   \big(\lambda , (a_1,\dots,a_n)\big) &\mapsto
%   \big(a_1,\dots,a_{i-1},\lambda a_i,a_{i+1},\dots,a_n\big) 
% \end{align*}
% These actions commute with each other, so we get an action $l$ of
% $(\C^*)^n$ on $\X(\sigma)$, given by
% \[ l\big((\l_1,\dots,\l_n),(a_1,\dots,a_n)\big) =
% \big(\l_1a_1,\dots,\l_na_n\big) \]

% Likewise, we can also multiply the rows of $A\in
% \X(\sigma)$ by non-zero complex numbers. This defines an action $l'$
% of $(\C^*)^m$ on $\X(\sigma)$. The two actions together satisfy the
% relation 
% \[ l\Big((\l,\dots,\l),l'\big((\l^{-1},\dots,\l^{-1}), A\big)\Big) =
% A \]
% since multiplying all rows by $\l^{-1}$ while multiplying all columns
% by $\l$ is the same as doing nothing.

% For the special case of $\X = \X(\Id)$, there is an additional
% symmetry that can be useful to know. By thinking of $A \in \X$ as a
% matrix, it is possible to transpose it. Checking the various
% determinants used in the definition of $\X$ show that they do not
% change after transposition. Hence transposing defines a homeomorphism
% $T : X_{m,n} \to X_{n,m}$, allowing us to switch the order of $m$
% and $n$.

% \section{The quotient space}

% % Rummene Y_{m,n}. 
% While the above definition of $\X$ gives an open subset of $\C^{mn}$,
% it has some issues that makes the space hard to
% work with. As we shall see shortly, the fundamental group and the
% cohomology ring of $\X$ both become larger as $n$ grows.
% This complicates things slightly, but it can be fixed by taking the
% quotient of a group action.
% \begin{definition}
%   The space $Y_{m,n}$ is the quotient space
%   \[ Y_{m,n} = \X / T^n \]
%   where $T^n$ acts on $X_{m,n}$ by scaling
%   the columns, as defined above:
%   \[ (\lambda_1,\dots,\lambda_n)\cdot (a_1,\dots,a_n) =
%   (\l_1a_1,\dots,\l_na_n) \]
% \end{definition}

% These spaces retain some of the structure described above. For
% example, the stabilization map $s$ is linear in $A$, which means it
% respects the group action.
% \[ s\big((\l_1,\dots,\l_n) \cdot A\big) = (\l_1,\dots,\l_n,1)\cdot
% s(A) \] 
% Hence $s$ descends to a stabilization map $s : Y_{m,n} \to
% Y_{m,n+1}$ and there is a direct limit,
% \[ Y_{m,\infty} = \varinjlim_{n} Y_{m,n} \]
% which can be thought of as $X_{m,\infty}/T^\infty$.

% The space $Y_{m,n}$ is homeomorphic to the subspace of $\X$ consisting
% of elements with the last $n$ coordinates of the determinant map
% $\Det$ equal to 1:
% \[ Y_{m,n} \cong \set{ A \in \X \delim\big\vert\delim \Det(A) \in
%   (\C^*)^{m-1}\times \set{1}^n }\]
% The group action gives an identification,
% \[ \X \cong Y_{m,n}\times (\C^*)^n \]
% This shows the claim about the fundamental group of $\X$, but it can
% also be used to get information about $Y_{m,n}$ from $\X$ or
% vice versa. For example, it allows us to calculate the cohomology of
% one when we know the cohomology of the other, by using the
% K\"unneth formula to add or remove the cohomology of $(\C^*)^n$.



%%% Local Variables: 
%%% mode: latex
%%% TeX-master: "main"
%%% End: 
 % Vis at rummene ikke er tomme (både
              % dimensions-argumentet og "stabilisering af de oplagte
              % elementer"-argumentet
              % skal nok være med. Kunne også beskrive
              % spektralfølgen.

% Hvad med dimension 2 tilfældet? Dette kunne inddrages fra Marcels
% noter og måske fylde et lille kapitel. Eventuelt før de generelle
% rum præsenteres.

% Hvor skal Coxeter-gruppe-halløjet ind? Muligvis allerede når man
% snakker om stabilisering. Og bevise sætningen om hvad man
% stabiliserer til.


% Dette skal vist næsten bare kopieres direkte over. Skift X til Y og
% brug måske lidt mere tid på at sætte det op for kvotientrummene også
% så man kan bruge det direkte på dem i stedet for at skulle igennem
% de ukvotienterede.
%% Hvis det ikke er gjort i tidligere kapitler, så lav udregningen af
%% H^1 generelt her et sted.


%%%%% Previous
\chapter{Cohomology calculations for $Y_{3,3}$}
\label{chapter:udregninger}

In this section we will calculate the cohomology of the space
$Y_{3,3}$, which is the set:
\[ Y_{3,3} = \set{
  \begin{pmatrix}
    a & b & g \\
    c & d & h \\
    e & f & i
  \end{pmatrix} \in \C^9 \sdel 
  \begin{matrix}
    a\neq 0,\, ad-bc \neq 0, \\
    adi+beh+cfg-deg-bci-afh \neq 0, \\
    di-fh \neq 0,\, i \neq 0
  \end{matrix} } \]

The way we will calculate this is by using the spectral sequence for
a filtration of $Y_{3,3}$. Some of the calculations have
been done in Sage and the relevant source code can be found at
\url{http://home.math.au.dk/sstoltze/del-a.html}. \fixme{Skift til github?}
% Some description of this space, properties, something?

\section{A filtration on $\Y$}

Before we start the calculation, we need some general setup. To work
with the space $Y_{m,n}$ and compute the cohomology, it will
be beneficial to build it up from simpler subspaces. This is done by
considering a filtration.

For $0 \leq p \leq m-1$, let $F_p$ be the subset
\[ F_p=\set{ (a_1,\dots,a_n) \in Y_{m,n} \mid a_n \text{ has at most } p
  \text{ zeroes}.} \]
$F_p$ is open in $\X$ since we stay inside the subspace when we
vary the coordinates of a point $A\in F_p$, as long as we do not
introduce new zeroes in the last column.
Together, the subspaces form an increasing sequence:
\[ F_0 \subset F_1 \subset \dots \subset F_{m-1} = Y_{m,n} \]

The general method of computing the homology and cohomology of
$Y_{m,n}$ is by using the spectral sequence of this
filtration, as described in \cite{hatcherss} and
\cite{mccleary}. Here we will focus on
cohomology, but similar arguments could of course be applied to
homology.

The spectral sequence has $E_1$-page given by
\[ E_1^{p,q} = H^{p+q}(F_p,F_{p-1}) \]
with differentials given by the boundary map in the long exact
sequence of the triple $(F_{p+1},F_p,F_{p-1})$:
\[ d_1 : H^{p+q}(F_p,F_{p-1}) = E_1^{p,q} \to E_1^{p+1,q} =
H^{p+q+1}(F_{p+1},F_p) \]

To use the spectral sequence, we need to compute the cohomology of the
pair $(F_p,F_{p-1})$. To calculate this we will use excision, so we
would like to find an open subset of $F_p$ that contains the
complement of $F_{p-1}$. For any choice of exactly $p$ 
different indices $\set{v_1,\dots,v_p}\subset \set{1,\dots,m-1}$,
define $Y_{\set{v_1,\dots,v_p}}$ to be the points $(a_1,\dots,a_n)\in
F_p$ which has $(a_n)_{v_i} = 0$ for all $i$. Since $\X$ is open, we
can get an open set in $F_p$ by inserting a small disc on these
zero-entries. More precisely, the set we get is
\[ Y_{\set{v_1,\dots,v_p}} \times D_p \cong \set{ A\in F_p \sdel 
  \begin{matrix} \displaystyle 
    \size{((a_n)_{v_1},\dots,(a_n)_{v_p})} <
    \min_{i\not\in \set{v_1,\dots,v_p}} \size{(a_n)_i}, \\
    \left(a_1,\dots,a_{n-1},a_n-t\sum_{i} (a_n)_{v_i}e_{v_i}\right)\in
    F_p \forall t \in [0,1]
  \end{matrix} } \]
This is an open set in $F_p$, with $Y_{\emptyset} = F_0$, and it
contains all points $(a_1,\dots,a_n)\in F_p$ where the entries of
$a_n$ indicated by $\set{v_1,\dots,v_p}$ are zero. For a different
choice of indices, the condition that the radius of the disc is
smaller than the norm of the non-zero entries in $a_n$ ensures
disjointness, so the intersection is empty:
\[ \left(Y_{\set{v_1,\dots,v_p}}\times D_p\right) \cap \left(
Y_{\set{u_1,\dots,u_p}}\times D_p\right) = \emptyset \quad \text{if }
\set{v_1,\dots,v_p} \neq \set{u_1,\dots,u_p} \]

The intersection of $\X\times D_p$ and $F_{p-1}$ corresponds to
removing 0 from the disc, since one of the entries inserted on the
last column has to be non-zero.
\[ \left(Y_{\set{v_1,\dots,v_p}}\times D_p \right) \cap F_{p-1} =
Y_{\set{v_1,\dots,v_p}} \times \left(D_p-0\right) \]
This means that the disjoint union of all of these form an open set in
$F_p$ with
\[ F_p = F_{p-1}\cup \left(\coprod_{\set{v_1,\dots,v_p}}
Y_{\set{v_1,\dots,v_p}}\times D_p\right) \]
and by excision, the inclusion of pairs
\[ \xymatrix{ \displaystyle \coprod_{\set{v_1,\dots,v_p}}
  Y_{\set{v_1,\dots,v_p}}\times (D_p,D_p-0) \ar@{^(->}[r] &
  (F_p,F_{p-1}) } \]
induces an isomorphism of cohomology groups,
\begin{align*}
  E_1^{p,q} &= H^{p+q}(F_p,F_{p-1}) \\
  &\cong H^{p+q}\left(\coprod_{\set{v_1,\dots,v_p}}
    Y_{\set{v_1,\dots,v_p}}\times (D_p,D_p-0)\right) \\   &\cong
  \bigoplus_{\set{v_1,\dots,v_p}}
  H^{p+q}(Y_{\set{v_1,\dots,v_p}}\times (D_p,D_p-0)) \\   &\cong
  \bigoplus_{\set{v_1,\dots,v_p}}
  H^{q-p}(Y_{\set{v_1,\dots,v_p}})\otimes H^{2p}(D_p,D_p-0) 
\end{align*}
The last isomorphism is given by the K\"unneth formula, and the fact
that the pair $(D_p,D_p-0)$ only has non-zero cohomology in degree
$2p$:
\[ H^q(D_p,D_p-0) \cong \widetilde H^{q-1}(D_p-0) = \widetilde
H^{q-1}(S^{2p-1}) =
\begin{cases}
  \Z & q = 2p \\
  0 & \text{otherwise}
\end{cases} \]


\section{The spectral sequence of $Y_{3,3}$}

We can apply the above filtration to get a spectral sequence of
$X=Y_{3,3}$. The groups we need to compute for the first page are:
\[ E_1^{p,q} \cong
\begin{cases}
  H^{q}(Y_\emptyset) & p=0 \\
  H^{q+1}(Y_{\set{1}} \times (D,D-0)) \oplus H^{q+1}(Y_{\set{2}} \times
  (D,D-0)) & p=1 \\
  H^{q+2}(Y_{\set{1,2}}\times (D_2,D_2-0)) & p = 2 \\
  0 & \text{otherwise}
\end{cases} \]
The four spaces that appear will be handled separately.

\subsection{The cohomology of $Y_{\set{1,2}}$}

Consider the space $Y_{\set{1,2}}$. This consists of the elements of
$Y_{3,3}$ that have two zeroes in the last column:
\[ Y_{\set{1,2}} = \set{
  \begin{pmatrix}
    a & b & 0 \\
    c & d & 0 \\
    e & f & i
  \end{pmatrix} \sdel
  \begin{matrix}
    a \neq 0, ad-bc \neq 0 \\
    adi-bci \neq 0 \\
    di \neq 0, i \neq 0
  \end{matrix}
} \]

The space is homeomorphic to $(\C^*)^3\times \C^2
\times \set{(b,c)\in\C^2 \delim \vert \delim bc\neq 1}$, with the
homeomorphism given by
\begin{align*}
  Y_{\set{1,2}} &\to (\C^*)^3\times \C^2
  \times \set{(b,c)\in\C^2 \delim \vert \delim bc\neq 1} \\
  \begin{pmatrix}
    a & b & 0 \\
    c & d & 0 \\
    e & f & i
  \end{pmatrix} &\mapsto \left(a,d,i,\frac{e}{i},\frac{f}{i},
    \frac{b}{a}, \frac{c}{d}\right)
\end{align*}

The space $\set{(b,c)\delim \vert\delim bc\neq 1}$ will be denoted
$Y$. To get the cohomology of this space, consider it as a subspace of
the one-point compactification $S^4 = \C^2\cup\set{\infty}$. Then the
complement is:
\[ S^4-Y = \set{(b,c)\in \C^2 \mid bc=1}\cup\set{\infty} \]
This can be identified with the wedge sum $S^1\vee S^2$ by considering
the subspace
\[ \C^* \cong \set{\left(b,\frac{1}{b} \right) \sdel b\in \C^*} =
(S^4-Y)-\set{\infty} \]
The point at infinity is glued to this space when either $b$ or $c$
is large, which for our copy of $\C^*$
corresponds to $b$ large or $b$ near 0. Hence what we get is
$S^2=\C\cup\set{\infty}$ with the identification $0 \sim \infty$. This
is homotopy equivalent to $S^1\vee S^2$.

Now we can apply Alexander duality to this space to compute
\[ \widetilde H_q(Y) \cong \widetilde H^{3-q}(S^4-Y) =
\begin{cases}
  \Z & q \in \set{1,2} \\
  0 & \text{otherwise}
\end{cases} \]
and hence
\[ H^q(Y) \cong
\begin{cases}
  \Z & q \in \set{0,1,2} \\
  0 & \text{otherwise}
\end{cases} \]

From all of this, we can calculate the cohomology of $Y_{\set{1,2}}$
\[ H^q(Y_{\set{1,2}}) =
\begin{cases}
  \Z & q = 0 \\
  \Z^4 & q = 1 \\
  \Z^7 & q = 2 \\
  \Z^7 & q = 3 \\
  \Z^4 & q = 4 \\
  \Z & q = 5 \\
  0 & \text{otherwise}
\end{cases} \]

Since we know the product structure from the K\"unneth formula, we can
name the generators and get the ring structure,
\[ H^*(Y_{\set{1,2}}) = \bigwedge\nolimits^*[a_{00},d_{00},i_{00}]\otimes
H^*(Y) \]
The two generators of $Y$ will be denoted $y_{00,1}$ for the degree
one generator and $y_{00,2}$ for the degree two generator. All other
generators have degree one. The ring
structure is the usual one on a tensor product, with the added
relation $y_{00,1}\cup y_{00,2} = 0$. The element $y_{00,1}\in H^1(Y)$
is given by the map
\[ Y\ni (b,c) \mapsto 1-bc\in \C^* \]
in cohomology, since there is another map
\[ \C^* \ni \l \mapsto (1,1-\l) \in Y \]
and the composition is the identity on $\C^*$. Applying $H^1$ shows
that the two maps give isomorphisms in degree 1.


\subsection{The cohomology of $Y_{\set{1}}$}

In $Y_{\set{1}}$ we have the first entry of the last column equal to
zero, which is the space
\[ Y_{\set{1}} = \set{
  \begin{pmatrix}
    a & b & 0 \\
    c & d & h \\
    e & f & i
  \end{pmatrix} \sdel
  \begin{matrix}
    a \neq 0, ad-bc \neq 0 \\
    adi + beh - bci - afh \neq 0 \\
    di -fh \neq 0, h \neq 0, i \neq 0 
  \end{matrix}
} \]

This space is homeomorphic to $(\C^*)^5 \times \C\times Y$:
\[ \begin{pmatrix}
  a & b & 0 \\
  c & d & h \\
  e & f & i
\end{pmatrix} \mapsto \left(
a, i, h, \frac{d}{h}-\frac{f}{i}, \frac{b}{h}-\frac{bc}{ah},
\frac{e}{i}, \frac{bhi}{adi-afh},\frac{c}{h}-\frac{e}{i}
\right) \] 
Since the cohomology of $Y$ is already known, applying the K\"unneth
formula gives
\[ H^q(Y_{\set{1}}) =
\begin{cases}
  \Z & q = 0 \\
  \Z^6 & q = 1 \\
  \Z^{16} & q = 2 \\
  \Z^{25} & q = 3 \\
  \Z^{25} & q = 4 \\
  \Z^{16} & q = 5 \\
  \Z^6 & q = 6 \\
  \Z & q = 7 \\
  0 & \text{otherwise}
\end{cases} \] 
As before, the ring structure is given by
\[ H^*(Y_{\set{1}}) =
\bigwedge\nolimits^*[a_{01},d_{01},f_{01},h_{01},i_{01}] \otimes H^*(Y)\]

\subsection{The cohomology of $Y_{\set{2}}$}

The space is
\[ Y_{\set{2}} = \set{
  \begin{pmatrix}
    a & b & g \\
    c & d & 0 \\
    e & f & i
  \end{pmatrix} \sdel
  \begin{matrix}
    a \neq 0, ad-bc \neq 0 \\
    adi +cfg -bci - deg \neq 0 \\
    di \neq 0, g\neq 0, i \neq 0 \end{matrix}
} \]
which can also be identified with $(\C^*)^5\times \C\times Y$ using
the homeomorphism
\[ \begin{pmatrix}
  a & b & g \\
  c & d & 0 \\
  e & f & i
\end{pmatrix} \mapsto \left(\frac{a}{g},d,i,g, 1+\frac{cfg}{adi}-
  \frac{eg}{ai}-\frac{bc}{ad}, \frac{f}{i},\frac{b}{a}, \frac{c}{d}
\right)\]
Since we already know the cohomology of this space, we know
\begin{align*}
  H^q(Y_{\set{2}}) = 
  \left(\bigwedge\nolimits^*[a_{10},d_{10},f_{10},g_{10},i_{10}]\otimes
  H^*(Y)\right)^q &=
  \begin{cases}
    \Z & q = 0 \\
    \Z^6 & q = 1 \\
    \Z^{16} & q = 2 \\
    \Z^{25} & q = 3 \\
    \Z^{25} & q = 4 \\
    \Z^{16} & q = 5 \\
    \Z^6 & q = 6 \\
    \Z & q = 7 \\
    0 & \text{otherwise}
  \end{cases}
\end{align*}

\subsection{The cohomology of $Y_{\emptyset}$}

This space is slightly more complicated than the others. The space
consists of the elements in $Y_{3,3}$ where all entries of the last
vector are non-zero.
\[ Y_{\emptyset} = \set{
  \begin{pmatrix}
    a & b & g \\
    c & d & h \\
    e & f & i
  \end{pmatrix} \sdel
  \begin{matrix}
    a \neq 0, ad-bc \neq 0, di-fh \neq 0, \\
    adi+beh+cfg-deg-bci-afh \neq 0, \\
    g\neq 0, h\neq 0, i\neq 0
  \end{matrix} }\]
The space is homeomorphic to $(\C^*)^7\times \set{ (x,y,z)\in \C^3
  \mid y-z-xyz = 1}$, with the homeomorphism given by
\begin{align*}
  \begin{pmatrix}
    a & b & g \\
    c & d & h \\
    e & f & i
  \end{pmatrix} \mapsto \bigg(&g, h, i, \frac{a}{g},
    \frac{f}{i}-\frac{d}{h},\\
  &D=\frac{ad-bc}{ah}, \\
  &C=\frac{c}{h}-\frac{a}{g}+\frac{cdi-deh}{fh^2-dhi} +
  \frac{beh-bci}{fgh-dgi},\\
  &\left.\frac{bC}{aD},\frac{c}{hC}-\frac{a}{gC},
    \frac{D}{C}\frac{eh-ci}{fh-di}\right)
\end{align*}

To get the cohomology, we need to compute the cohomology of the space:
\[ Z=\set{(x,y,z)\in\C^3 \mid y-z-xyz = 1} \]
This is the same as the cohomology of $S^2$, but this is hard to see
directly. Instead consider the space
\[ \widetilde Z = \set{(b,c,d,e) \in \C^4 \mid d\neq
  bc, c\neq de} \]
This space is homeomorphic to $Z\times (\C^*)^2$:
\begin{align*}
  (b,c,d,e) \mapsto \left(
    \frac{b(c-de)}{d-bc}, \frac{c}{c-de}, \frac{e(d-bc)}{c-de},
  d-bc, c-de\right)
\end{align*}
We will use Alexander duality to compute the cohomology as we did for
$Y$. Consider the complement in $S^8$,
\[ S^8-\widetilde Z = \set{(b,c,d,e) \mid d = bc} \cup \set{(b,c,d,e)
  \mid c = de} \cup \set{\infty} \]
and denote the two sets as $W_1$ and $W_2$. Then
\begin{align*}
  W_1\cap W_2 &= \set{(b,c,d,e) \mid d=bc,c=de}\cup\set{\infty}
\end{align*}
We glue the point at infinity to this space when $c$ is large, but
this is the same as either $d$ or $e$ being large. Hence there is a
homeomorphism,
\begin{align*}
  W_1\cap W_2 &\cong \set{(b,d,e) \mid d=bde}\cup\set{\infty} \\
  &= \set{(b,d,e)\mid 0 = (be-1)d}\cup\set{\infty} \\
  &= \set{(b,d,e) \mid d = 0} \cup \set{(b,d,e) \mid
    be=1}\cup\set{\infty}
\end{align*}

These new spaces have computable homology, since they can be
identified with spaces built from spheres. The first is relatively
easy,
\[ Z_1 = \set{(b,d,e)\mid d = 0}\cup \set{\infty} =
\set{(b,0,0,d)}\cup\set{\infty} \cong S^4 \]
The second is slightly harder.
\[ Z_2 =\set{(b,d,e)\mid be=1}\cup\set{\infty}=
\set{\left(b,\frac{d}{b},d,\frac{1}{b} \right) }\cup\set{\infty} \]
This set is homeomorphic to the one-point compactification of
$\set{(b,d)\in \C^*\times \C}$. By \cite[Chapter~2.3]{dupontk}, this
is given by the smash product of the one-point compactifications,
\[ Z_2\simeq (\C^*\times\C)^+ \cong (\C^*)^+\wedge\C^+ \cong (S^1\vee
S^2) \wedge S^2 = S^3 \vee S^4 \]
The intersection of these spaces is
\[ Z_1\cap Z_2 = \set{\left(b,0,0,\frac{1}{b}\right)}\cup\set{\infty}
\simeq S^1\vee S^2 \]

From all of this, we get a commutative diagram of inclusions,
\[ \xymatrix{
  S^1\vee S^2 \ar[r]\ar[d] & S^3\vee S^4 \ar[d] \\
  S^4 \ar[r] & W_1\cap W_2
} \]
The two maps from $S^1\vee S^2$ are zero in cohomology (when the
degree is greater than zero) since in each
degree either the domain or the codomain is zero. By applying
Mayer-Vietoris, we get
\[ H^q(W_1\cap W_2) =
\begin{cases}
  \Z & q=0 \\
  0 & q = 1 \\
  \Z & q = 2 \\
  \Z^2 & q = 3 \\
  \Z^2 & q =4 \\
  0 & \text{otherwise}
\end{cases} \]
where the case $q=0$ follows from $W_1\cap W_2$ being connected
and $q = 1$ follows from writing out the relevant part of the
long exact sequence,
\[ \xymatrix{ H^0(S^3\vee S^4)\oplus H^0(S^4) \ar@{->>}[r] &
H^0(S^1\vee S^4) \ar[r] & H^1(W_1\cap W_2) \ar[r] & 0} \]

$W_1$ and $W_2$ are both a copy of $S^6$ since we can get one of
the coordinates from two of the others. Hence we get another diagram,
\[ \xymatrix{
  W_1\cap W_2 \ar[r]\ar[d] & S^6 \ar[d] \\
  S^6 \ar[r] & S^8-\widetilde Z
} \]
and again the map from the intersection is zero. Repeating the above,
we can compute:
\[ H^q(S^8-\widetilde Z) =
\begin{cases}
  \Z & q = 0 \\
  0 & q = 1 \\
  0 & q = 2 \\
  \Z & q = 3 \\
  \Z^2 & q = 4 \\
  \Z^2 & q = 5 \\
  \Z^2 & q =6 \\
  0 & \text{otherwise}
\end{cases} \]
Applying Alexander duality then gives
\[ H^q(\widetilde Z) =
\begin{cases}
  \Z & q = 0 \\
  \Z^2 & q = 1 \\
  \Z^2 & q = 2 \\
  \Z^2 & q = 3 \\
  \Z & q = 4 \\
  0 & \text{otherwise}
\end{cases} \]
from which we get that the cohomology of $Z$ is the same as the
cohomology of $S^2$ by factoring out the torus. The generator of
$H^2(Z)$ will be called $z_{11}$ with $z_{11}^2 = 0$.

All in all, this shows
\[ H^q(Y_{\emptyset}) =
  \begin{cases}
    \Z & q = 0 \\
    \Z^7 & q = 1 \\
    \Z^{22} & q = 2 \\
    \Z^{42} & q = 3 \\
    \Z^{56} & q = 4 \\
    \Z^{56} & q = 5 \\
    \Z^{42} & q = 6 \\
    \Z^{22} & q = 7 \\
    \Z^{7} & q = 8 \\
    \Z & q = 9 \\
    0 & \text{otherwise}
  \end{cases} \] 
and the generators are
\[ H^*(Y_{\emptyset}) =
\bigwedge\nolimits^*[a_{11},d_{11},g_{11},h_{11},i_{11},C_{11},D_{11}]
\otimes H^*(Z) \]

\subsubsection{A note on the cohomology of $Z$}

The space $Z$ defined above is a manifold, which we will prove by
applying the implicit function theorem. Consider the formula defining
$Z$,
\[ (x,y,z) \mapsto y-z-xyz-1 \]
Taking partial derivatives gives the three functions
\begin{align*}
  (x,y,z) &\mapsto yz \\
  (x,y,z) &\mapsto 1-xz \\
  (x,y,z) &\mapsto -1-xy
\end{align*}
If the first function is zero, then at least one of the two others
must be non-zero. By the implicit function theorem, the coordinate
that has non-zero derivative can be written as a smooth function of
the other coordinates. This shows that $Z$ is locally the graph of a
smooth function, so it is a smooth manifold.

Later on it will be necessary to have a description of $H^2(Z)$. This
is done by considering the subspace
\[ R = \set{(x,y,z) \in \R^3 \sdel x > 0, y < 0, xy < -1, z =
  \frac{y-1}{1+xy}>0 } \subset Z \]
This space is closed in $Z$ since any sequence in $R$ that converges
in $Z$ must also converge in $R$.
Consider a subbundle $NR \cong R\times D$ of the normal bundle of $R$
in the tangent bundle $TZ$, such that the exponential map is a
diffeomorphism on $NR$. This gives us a diagram
\[ \xymatrix{ H^2(NR,NR_0) & H^2(Z,Z-R) \ar[l]_\cong \ar[r]^>>>>>{j^*} &
  H^2(Z) } \]
where $NR_0$ is $NR$ without the zero section, $NR_0 \cong
R\times(D-0)$. The isomorphism is given by excision, since $Z = NR
\cup (Z-R)$. The space $R$ is homeomorphic to $\R^2$, so we have 
\[ H^2(Z,Z-R) \cong H^2(NR,NR_0) \cong H^2(R\times (D,D-0)) \cong
\Z \] 
From the diagram, $j^*(1)$ gives us an element of $H^2(Z)$. To see
that this is a generator, consider the torus 
\[ \widehat T = \set{(x,y,z)\in Z \mid \size{x} = \size{y} = 2 }
\subset Z \]
This intersects $R$ in exactly one point, namely $p=(2,-2,1)$. By
taking the intersection with $\widehat T$ we can extend the previous
diagram,
\[ \xymatrix{ H^2(NR,NR_0) \ar[d]^{i^*} & H^2(Z,Z-R) \ar[d]^{i^*}
  \ar[l]_<<<<<\cong \ar[r]^>>>>>{j^*} & H^2(Z) \ar[d]^{i^*} \\
  H^2(p\times (D,D-0)) & H^2(\widehat T,\widehat T-p)
  \ar[r]^>>>>>{\cong} \ar[l]_>>>>>\cong & H^2(\widehat T) 
} \]
The left $i^*$ is an isomorphism since the cohomology in degree two is
given by the pair $(D,D-0)$, which does not change under the
inclusion. This shows that the inclusion $\widehat T \to Z$ is an
isomorphism in cohomology, so $j^*(1)$ is a generator of
$H^2(Z)$. Taking the same diagram with cohomology replaced with
homology reverses all arrows but is otherwise exactly the same. This
shows that to calculate the homology class of a cycle $\sigma$ in
$H_2(Z)$, we can count the number of points of intersection between
$\sigma$ and $R$ with orientation, since each of these will contribute
either $1$ or $-1$, depending on orientation, to the homology class in
$(Z,Z-R)$.

There is a smiliar argument with the space $Y$ from before and the
subspace $R_Y = \set{(b,c) \in \R^2 \mid b > 0, c > 0, bc > 1}$. In
particular, the generator of $H_2(Y)$ can be taken as
$i_*\left([\widehat T]\right)$, where $[\widehat
  T]$ is a generator of $H_2(\widehat T)$ for the torus $\widehat T =
  \set{(b,c)\mid \size{b}=\size{c} = 2}$.

\subsection{The first page}

With all of the above, we can draw the $E_1$-page of the spectral
sequence. Remember that this is
\[ E_1^{p,q} \cong
\begin{cases}
  H^{q}(Y_\emptyset) & p=0 \\
  H^{q+1}(Y_{\set{1}} \times (D,D-0)) \oplus H^{q+1}(Y_{\set{2}} \times
  (D,D-0)) & p=1 \\
  H^{q+2}(Y_{\set{1,2}}\times (D_2,D_2-0)) & p = 2 \\
  0 & \text{otherwise}
\end{cases} \]

By inserting the groups that were calculated above into this, we get
Figure \ref{e1} with the non-zero differentials moving directly to
the right. To get the $E_2$-page of the spectral sequence, we need to
compute the homology of the differentials.

\begin{figure}[ht]
  \[ \xy
  (5,4); (5,-80); **\dir{-}; (5,4); *\dir{>};
  (5,-80); (42,-80); **\dir{-}; *\dir{<};
  (0,-82)*{q}; (5,-84)*{p};
  (1,-84); (4,-81); **\dir{-};
  (0,0)\xymatrix@R=1em@C=1.7em{
    9 & \Z & 0 & 0 \\
    8 & \Z^7 \ar[r]& \Z^2 & 0 \\
    7 & \Z^{22} \ar[r]& \Z^{12} \ar[r]& \Z \\
    6 & \Z^{42} \ar[r]& \Z^{32} \ar[r]& \Z^{4} \\
    5 & \Z^{56} \ar[r]& \Z^{50} \ar[r]& \Z^{7} \\
    4 & \Z^{56} \ar[r]& \Z^{50} \ar[r]& \Z^{7} \\
    3 & \Z^{42} \ar[r]& \Z^{32} \ar[r]& \Z^{4} \\
    2 & \Z^{22} \ar[r]& \Z^{12} \ar[r]& \Z \\
    1 & \Z^{7} \ar[r]& \Z^{2} & 0 \\
    0 & \Z & 0 & 0 \\
    & 0 & 1 & 2
  } \endxy \]
  \caption{The $E_1$-page with differentials.}
  \label{e1}
\end{figure}

\section{Differentials on $E_1$}
\label{sec:orientation}

By definition, the differentials are the boundary maps of the long
exact sequence of the triple $(F_{p+1},F_p,F_{p-1})$,
\[ d_1 : E_1^{p,q} = H^{p+q}(F_p,F_{p-1}) \to H^{p+q+1}(F_{p+1},F_p) =
E_1^{p+1,q} \]
We will consider the two cases $p=0$ and $p=1$ separately.

\subsection{First column}

For $p = 0$, we are looking at the boundary map of the pair
$(F_1,F_0)$. Consider the inclusion of pairs,
\[ i : Y_{\set{1}}\times (D,D-0) \sqcup Y_{\set{2}}\times (D,D-0) \to
(F_1,F_0) \]
This is an isomorphism in cohomology and gives a map between the long
exact sequences that makes the following diagram
commute:
\[ \xymatrix{
H^{p+q}(F_0) \ar[r]^>>>>>>>>>>>>>>>>>>>>{d_1} \ar[d]^{i^*} &
H^{p+q+1}(F_1,F_0) \ar[d]^{i^*}_{\cong} \\
H^{p+q}((Y_{\set{1}}\sqcup Y_{\set{2}})\times (D-0))
\ar[r]^>>>>>{\delta} & H^{p+q+1}((Y_{\set{1}}\sqcup Y_{\set{2}})\times
(D,D-0)) } \]
The bottom map is the boundary map from the long exact sequence of the
pair $(D,D-0)$. This map is well-known and will be treated at the end
of this section, so the only missing part is to calculate the
inclusion $i^*$. Since we are working with a disjoint union, we only
need to calculate the inclusion on each factor.

\subsubsection{The inclusion $Y_{\set{1}}\times (D-0) 
  \hookrightarrow Y_\emptyset$}

We can write the inclusion out in coordinates as
\begin{align*}
  (\C^*)^5\times Y \times (D-0) \longrightarrow
  (&\C^*)^7\times Z \\
  (a,d,f,h,i,b,c,g) \mapsto \bigg(&g,h,i,\frac{a}{g},-f,d,\\
  & C =
  c-\frac{a}{g}-\frac{cd}{f}-bc^2+\frac{abc}{g}, \\
  & \frac{bfC}{d},\frac{c}{C}-\frac{a}{g C},-\frac{cd}{fC}\bigg) 
\end{align*}

On the generators, this is given by
\begin{align*}
  i^* : H^*(Y_{\emptyset}) &\to H^*(Y_{\set{1}}\times (D-0)) \\
  a_{11} & \mapsto a_{01} - g_{01} \\
  d_{11} &\mapsto f_{01} \\
  g_{11} &\mapsto g_{01} \\
  h_{11} &\mapsto h_{01} \\
  i_{11} &\mapsto i_{01} \\
  D_{11} &\mapsto d_{01} \\
  C_{11} &\mapsto a_{01} - g_{01} + y_{01,1} \\
  z_{11} &\mapsto y_{01,2} + (a_{01}-d_{01}+f_{01}-g_{01})y_{01,1}
\end{align*}
where $g_{01}$ is the generator of $H^1(D-0)$.

In degree 1, this is calculated by dualizing to homology, writing out
the map on generators and working
out what they map to. For $z_{11}$ it is slightly
harder. To calculate this, we use the submanifold $R \subset Z$
defined previously as
\[ R = \set{(x,y,z)\in \R^3 \delim \Big\vert\delim x > 0, y < 0, z =
  \frac{y-1}{1+xy}>0} \]
by choosing a generator $\alpha = [\sigma]$ of $H_2(Y_{\set{1}}\times
(D-0))$ and calculating the intersection of the image of $f = \pr_Z
\circ i\circ \sigma$ and $R$. Note that if the coordinates of $Y$ are
fixed at $(0,1)$, then the image of
\[ f(a,d,f,h,i,0,1,g) =
\left(0,\frac{1}{C}-\frac{a}{gC},-\frac{d}{fC} \right)\]
never intersects $R$ since $x = 0$. Likewise, the coordinates $i$ and
$h$ do not appear in the projection to $Z$ and can not influence the
result. Hence we only have to check the duals of the following
generators in degree 2:
\[ a_{01} y_{01,1}, d_{01} y_{01,1}, f_{01} y_{01,1}, y_{01,1} g_{01},
y_{01,2} \]
These will be handled individually. All non-relevant coordinates will
be set to 1, except $g$ which will be denoted $\e$ to make it clear
that this coordinate is small. If $y_{01,1}$ is involved, we choose the
generator
\[ \sigma(\lambda) = (1-\lambda,1) \qquad \lambda\in S^1 \]
as this gives
\[ y_{01,1} = [t \mapsto 1-(1-\lambda)\cdot 1] = [t\mapsto \lambda] \]

\paragraph{$\mathbf{a_{01} y_{01,1}}$:}
The relevant map is
\[ (a,\lambda) \mapsto \left( (1-\l)C, \frac{1}{C}-\frac{a}{\e
    C},\frac{1}{C}\right)\]
where $C = -1-\frac{a\l}{\e}$. If this intersects $R$ then $C$ has
to be real and positive by looking at the last coordinate, which
forces $\l$ real and negative by the first coordinate and $a$ real and
positive by the second. Hence there is at most one intersection, at
$(1,-1)$:
\[ f(1,-1) = \left( \frac{2-2\e}{\e},-1,\frac{\e}{1-\e}\right) \]
By inspection, this satisfies the conditions and hence belongs to
$R$, so we conclude that the image intersects once. This gives the
following formula:
\[ \langle f^*(z_{11}), a_{01}y_{01,1} \rangle = \langle
z_{11},f_*(a_{01}y_{01,1})\rangle = \pm 1 \]
To check the sign, we need to check if this map preserves the
orientation of the torus in $R$. This is done by differentiating the
map to get a map of tangent spaces
\[ Tf : TT^2 \to TR \]
and checking that the determinant of this map is
positive. This has been done in Sage, see Appendix
\ref{sec:sageorientation} for the code.

\paragraph{$\mathbf{d_{01} y_{01,1}}$:}
In this case the map becomes
\[ (d,\l) \mapsto \left( \frac{(1-\l)C}{d},\frac{\e-1}{\e C},
  \frac{d}{C} \right) \]
and by a similar analysis, this intersects in exactly the
point $f(1,-1)$, which is the same as before. However, this time the
orientation is reversed.

\paragraph{$\mathbf{f_{01} y_{01,1}}$:}
This has the same point of intersection and preserves the orientation.

\paragraph{$\mathbf{y_{01,1}g_{01}}$:}
This also intersects and preserves the orientation.

\paragraph{$\mathbf{y_{01,2}}$:}
This is calculated by the intersection points of the fundamental class
of the torus in $Y$, which was mentioned earlier:
\[ \widehat T = \set{(b,c) \in \C^2 \mid \size{b}=\size{c} = 2}
\subset Y \]
The map becomes
\[ (b,c) \mapsto \left( bC, \frac{c}{C}-\frac{1}{\e C},\frac{c}{C}
\right) \]
with $C = -bc^2-\frac{1-bc}{\e}$. We will write $x$ for the first
coordinate, $y$ for the second and $z$ for the third. An analysis as
before gives $bc = xz > 0$ and $c = \frac{1}{\e(z-y)} \in
\R$. So both $b$ and $c$ must be real and $c$ should be positive. Then
$xy = bc - b\e^{-1}<0$ shows that $b$ must be positive since $bc$ is
positive. The only possible point of intersection is
\[ f(2,2) = \left( \frac{6-16\e}{\e},
  \frac{2\e-1}{3-8\e},\frac{2\e}{3-8\e} \right) \]
Since $\e$ is very small, this intersects $R$ and the orientation is
preserved.

\subsubsection{The inclusion $Y_{\set{2}}\times (D-0) 
  \hookrightarrow Y_\emptyset$}

The inclusion in coordinates is
\begin{align*}
  (\C^*)^5\times Y \times (D-0) \longrightarrow (&\C^*)^7\times Z \\
  (a,d,f,g,i,b,c,h) \mapsto 
  \bigg(g&,h,i,a,-\frac{d}{h},\\
  & D=\frac{d}{h}(1-bc), \\
  & C= -af-\frac{a^2bh}{d}+\frac{a^2bhf}{d}+\frac{a^2b^2ch}{d}, \\
  & \frac{bC}{D}, \frac{cd}{hC}-\frac{a}{C},
  \frac{cD}{C}-\frac{ahD}{dC}+\frac{afhD}{dC}+
  \frac{abchD}{dC}\bigg)
\end{align*}

We calculate as before and get the following on generators:
\begin{align*}
  i^* : H^*(Y_{\emptyset}) &\to H^*(Y_{\set{2}}\times (D-0)) \\
  a_{11} & \mapsto a_{10} \\
  d_{11} &\mapsto d_{10} - h_{10} \\
  g_{11} &\mapsto g_{10} \\
  h_{11} &\mapsto h_{10} \\
  i_{11} &\mapsto i_{10} \\
  D_{11} &\mapsto d_{10} + y_{10,1} - h_{10} \\
  C_{11} &\mapsto f_{10} + a_{10} \\
  z_{11} &\mapsto y_{10,2} + (a_{10}-d_{10}+f_{10}+h_{10})y_{10,1}
\end{align*}
where $h_{10}$ is the generator of $H^1(D-0)$.

The calculations have been left out, but are very similar to the
previous ones.

\subsection{Second column}

The differential on the second column is the boundary map
\[ \delta : H^q((Y_{\set{1}}\sqcup Y_{\set{2}})\times (D,D-0)) \cong
H^q(F_1,F_0) \to H^{q+1}(F_2,F_1) \] 

To calculate this map, consider the inclusion of triples
\[ Y_{\set{1,2}}\times (D_2,D_2-0,(D-0)\times (D-0)) \hookrightarrow
(F_2,F_1,F_0) \]
given by inserting a small disc, sphere or torus on the zero-entries
of $Y_{\set{1,2}}$. There are homeomorphisms of pairs,
\begin{align*}
  \big(D_2-0,(D-0)\times(D-0)\big) \cong \big(D\times (D-0) \sqcup
  (D-0)\times D, (D-0)\times(D-0)\big) \\
  \big(D_2,D_2-0\big) \cong \big(D\times D,D\times
  (D-0)\cup (D-0)\times D\big) =  (D,D-0)\times (D,D-0)
\end{align*}
These gives a commuting diagram,
\[ \xymatrix{ H^q(F_1,F_0) \ar[r] \ar[d] & H^{q+1}(F_2,F_1)
  \ar[d]^\cong \\ 
  H^q(Y_{\set{1,2}}\times (D_2-0,(D-0)\times (D-0))) \ar[r]
%  \ar@{=}[d] 
&
  H^{q+1}(Y_{\set{1,2}} \times (D_2,D_2-0)) %\ar@{=}[d] 
%\\
%  H^q(Y_{\set{1,2}}\times ((D,D-0)\times (D-0) \sqcup (D-0)\times
%  (D,D-0))) \ar[r] & H^{q+1}(Y_{\set{1,2}}\times (D,D-0)\times (D,D-0))
} \]
The marked isomorphism follows from excision and the left map is
induced by an inclusion. If we apply the
homeomorphism of pairs to the bottom of the diagram, the boundary map
at the bottom becomes the sum of the two boundary maps,
\[ \xymatrix@R=0.1em{ H^q(Y_{\set{1,2}}\times (D,D-0)\times (D-0))
  \ar[]!<15ex,0ex>;[rd]!<-20ex,0ex>^{\delta} \\
  & H^{q+1}(Y_{\set{1,2}}\times (D,D-0)\times (D,D-0)) \\
  H^q(Y_{\set{1,2}}\times (D-0)\times (D,D-0))
  \ar[]!<15ex,0ex>;[ur]!<-20ex,0ex>_\delta } \]
The boundary maps are given by the long exact sequence of the pair
$(D,D-0)$, so we only need to consider the inclusion on the left. By
choosing the discs sufficiently small and using the identification of
$H^q(F_1,F_0)$ defined in the first section, 
we see that the inclusion maps the space $Y_{\set{1,2}}\times
(D,D-0)\times (D-0)$ to $Y_{\set{1}}\times(D,D-0)$ and
$Y_{\set{1,2}}\times (D-0)\times (D,D-0)$ to $Y_{\set{2}}\times
(D,D-0)$. To get the differential, we must calculate the sum of these
two maps:
\begin{align*}
  i_1^* : H^q(Y_{\set{1}}\times (D,D-0)) &\to H^q(Y_{\set{1,2}}\times
  (D,D-0)\times (D-0)) \\
  i_2^* : H^q(Y_{\set{2}}\times (D,D-0)) &\to H^q(Y_{\set{1,2}}\times
  (D-0) \times (D,D-0))
\end{align*}

\subsubsection{The inclusion $Y_{\set{1,2}}\times (D,D-0) 
\times (D-0) \to Y_{\set{1}} \times (D,D-0)$}   

In coordinates, the map becomes
\begin{align*}
  (\C^*)^3 \times Y \times (D,D-0)\times (D-0) &\to (\C^*)^5 \times Y
  \times (D,D-0) \\
  (a,d,i,b,c,g,h) &\mapsto \left(a,i,h,\frac{d}{h},
  \frac{d}{h}(1-bc), \frac{bh}{d},\frac{dc}{h},g\right)
\end{align*}

Again, we calculate the map on cohomology by dualizing to homology
and working with the dual generators. To check what happens to the
elements of $H^q(Y)$, we count intersections with the subspace $R_Y =
\set{(b,c)\in \R^2 \mid b > 0, c > 0, bc > 1}$ of $Y$, exactly as for
$Z$ above.
\begin{align*}
  i_1^* : H^q(Y_{\set{1}}\times (D,D-0)) &\to H^q(Y_{\set{1,2}}\times
  (D,D-0)\times (D-0)) \\
  a_{01} &\mapsto a_{00} \\
  d_{01} &\mapsto d_{00}-h_{00} + y_{00,1} \\
  f_{01} &\mapsto d_{00} - h_{00} \\
  g_{01} &\mapsto g_{00} \\
  h_{01} &\mapsto h_{00} \\
  i_{01} &\mapsto i_{00} \\
  y_{01,1} &\mapsto y_{00,1} \\
  y_{01,2} &\mapsto y_{00,2} - d_{00}y_{00,1} + h_{00}y_{00,1}
\end{align*}
We denote the generator of $H^2(D,D-0)$ in
$H^*(Y_{\set{1}}\times (D,D-0))$ by $g_{01}$ and for
$H^*(Y_{\emptyset} \times (D,D-0) \times (D-0))$ we consider $g_{00}\in
H^2(D,D-0)$ and $h_{00} \in H^1(D-0)$. 

\subsubsection{The inclusion $Y_{\set{1,2}}\times (D-0)
\times (D,D-0) \to Y_{\set{2}}\times (D,D-0)$}
This is exactly as above for the map
\begin{align*}
  (\C^*)^3 \times Y \times (D-0)\times (D,D-0) &\to (\C^*)^5 \times
  Y\times (D,D-0) \\ 
  (a,d,i,b,c,g,h) &\mapsto
  \left(\frac{a}{g},d,i,g,1-bc,b,c,h\right)
\end{align*}
In cohomology, this becomes
\begin{align*}
  i_2^* : H^q(Y_{\set{2}}\times (D,D-0)) &\to H^q(Y_{\set{1,2}}\times
  (D-0) \times (D,D-0)) \\
  a_{10} &\mapsto a_{00} - g_{00} \\
  d_{10} &\mapsto d_{00} \\
  f_{10} &\mapsto y_{00,1}\\
  g_{10} &\mapsto g_{00} \\
  h_{10} &\mapsto h_{00} \\
  i_{10} & \mapsto i_{00} \\
  y_{10,1} &\mapsto y_{00,1} \\
  y_{10,2} &\mapsto y_{00,2}
\end{align*}
with $g_{01}\in H^2(D,D-0)$ inside $H^*(Y_{\set{2}}\times (D,D-0))$ and
$g_{00}\in H^2(D,D-0)$, $h_{00}\in H^1(D-0)$ inside
$H^*(Y_{\emptyset}\times(D-0)\times(D,D-0))$.

\subsubsection{Boundary maps}
\label{disk}
We have now reduced everything to calculating some well-known boundary
maps for the pair $(D,D-0)$. By writing out the long exact sequence,
the boundary map
\[ \delta : H^*(D-0) \to H^*(D,D-0) \]
is an isomorphism in degree 1,
\[ H^1(D-0) \cong H^2(D,D-0) \]
and zero in all other degrees. The boundary map for the pair $X\times
(D,D-0)$ is $\Id\otimes \delta$ in tensor-notation. So the boundary
map in the cases above picks out the ``new'' coordinate added to
$X\times (D-0)$ and throws away everything that does not depend on
it. To be more precise, the map is given on generators by:
\begin{align*}
  H^*(X)\otimes H^*(D-0)) &\to H^*(X)\otimes H^*(D,D-0)) \\
  x\otimes d &\mapsto
  \begin{cases}
    x \otimes \delta d & d\in H^1(D-0) \\
    0 & \text{otherwise}
  \end{cases}
\end{align*}

\section{Turning the page}
\label{sec:rand}

We can now use the above to calculate the $E_2$-page. We have
calculated the groups appearing on the first page and the
differential. To get the second page, we need to compute the homology
groups of $d_1$. This reduces to linear algebra over the integers, and
is done using Sage. The result can be found in Figure \ref{e2}.

\begin{figure}[ht]
  \[ \xy
  (5,4); (5,-80); **\dir{-}; (5,4); *\dir{>};
  (5,-80); (42,-80); **\dir{-}; *\dir{<};
  (0,-82)*{q}; (5,-84)*{p};
  (1,-84); (4,-81); **\dir{-};
  (0,0)\xymatrix@R=1em@C=1.7em{
    9 & \Z & 0 & 0 \\
    8 & \Z^5 & 0 & 0 \\
    7 & \Z^{11} & 0 & 0 \\
    6 & \Z^{15} & \Z & 0 \\
    5 & \Z^{16} & \Z^{3} & 0 \\
    4 & \Z^{16} & \Z^{3} & 0 \\
    3 & \Z^{15} & \Z & 0 \\
    2 & \Z^{11} & 0 & 0 \\
    1 & \Z^{5} & 0 & 0 \\
    0 & \Z & 0 & 0 \\
    & 0 & 1 & 2
  } \endxy \]
  \caption{The $E_2$-page.}
  \label{e2}
\end{figure}

The differentials from $E_2$ onwards are given as maps:
\[ d_r : E_r^{p,q} \to E_r^{p+r,q+1-r} \]
For $r\geq 2$, this increases $p$ by at least 2. This means that
the domain or codomain of $d_r$ will always be zero, so taking
homology does not change anything. Hence $E_2 = E_\infty$, the
spectral sequence collapses at the 2\textsuperscript{nd} term, and
there are no extension problems since all the groups are free and
abelian. This allows us to read off the cohomology of $Y_{3,3}$. Since
we are also interested in factoring out the $T^3$ torus, we note this
as well.
\[ H^q(Y_{3,3}) = E_2^{0,q} \oplus E_2^{1,q-1} = 
\begin{cases}
  \Z & q = 0 \\
  \Z^{5} & q = 1\\
  \Z^{11} & q = 2\\
  \Z^{15} & q = 3\\
  \Z^{17} & q = 4\\
  \Z^{19} & q = 5\\
  \Z^{18} & q = 6\\
  \Z^{12} & q = 7\\
  \Z^{5} & q = 8\\
  \Z & q = 9\\
  0 & \text{otherwise} 
\end{cases} \qquad
H^q(X_{3,3}) =
\begin{cases}
  \Z & q = 0 \\
  \Z^{2} & q=1 \\
  \Z^{2} & q=2 \\
  \Z^{2} & q=3 \\
  \Z^{3} & q=4 \\
  \Z^{2} & q=5 \\
  \Z & q=6 \\
  0 & \text{otherwise}
\end{cases}
\]
We can also read off the ring structure from the spectral sequence, as
it is induced from the spaces we start with. Since we do not
need it, it will not be included here.
% Også noget med \pi_1? Eg. surjektivitet eller lignende?


%%% Local Variables: 
%%% mode: latex
%%% TeX-master: "main"
%%% End: 
 % Skal nok tages nogenlunde direkte fra
                              % del A-rapporten. Måske kan man smide
                              % udregningen af H^1 ind her?


% Løkker!
%% Hvis løkkerummet ikke er blevet præsenteret endnu, så start med
%% dette.
%% Hvis sætningerne fra Milnor ikke er blevet introduceret, så lav et
%% kort afsnit om dem her.

% Sætningen
%% Med bevis.
%% Variationer: \R og \H. Med bevis (måske mere med skitse hvor man
%% overvejer hvor \H er forskellig fra \C. \R er det oplagte tilfælde
%% at udelade for ikke at gøre det hele ensformigt. Måske en kort note
%% om hvor sammenhæng af skalarerne spiller ind.

% Stabiliserings-argumenter? 
%% Måske kan det Marcel laver smides ind her for generelt m, ellers
%% skal det gemmes til senere.

%%%%% Previous
\chapter{Loop spaces} 
\fixme{Skal stadig nok revideres en del.}

We now return to the general case. We want to show that the limit
space
$X_{m,\infty}$ is equivalent to the space of loops in the flag
manifold $\GLB{m}$.

\begin{definition}
  The loop space $\Omega = \Omega(\GLB{m})$ is the space
  \[ \Omega = \set{\g : [0,1] \to \GLB{m} \mid \g(0) = \g(1) =
    [\Id]} \]
  with the compact-open topology. 
\end{definition}

\begin{remark} \fixme{Flyt til introduktions-området. Hvis intet
    introduktions-område, så skriv om og inkluder de vigtigste
    resultater fra Milnor her uden bevis}
  Following \cite[Chapter 17]{milnor}, the compact-open topology is
  also the topology given by two loops being close if they are close
  at every point, with the metric
  \[ d(\g,\g') = \max_{t\in[0,1]} \widetilde d (\g(t),\g'(t)), \]
  where $\widetilde d$ is a metric on $\GLB{m}$ induced from a
  Riemannian metric. The space $\Omega$ is homotopy equivalent to the
  space of piecewise smooth loops with a slightly modified metric, and
  we will generally choose the most convenient model for our
  arguments.
\end{remark}

\section{The map $X_{m,n}\to\Omega$}

As has been done before, we will start by considering $Y_{m,n}$ and
then quotient out with the torus $\T^n$ afterwards.
For $A = [a_1,\dots,a_n]\in Y_{m,n}$ we
can define a path $ \g_A$ in $\GL_m$ by starting with
$\g_A(0) = \Id$. Since the matrix $[e_2,\dots,e_m,a_1]$ is
invertible by the definition of $Y_{m,n}$, we know that $a_1$ is linearly
independent of $e_2,\dots,e_m$. Hence it can be written as
\[ a_1 = \sum_{i=1}^m c_i e_i, \]
with $c_1 \neq 0$. Rewriting the expression leads to the formula
\[  \frac{1}{c_1} a_1= e_1 + \frac{1}{c_1} \left(\sum_{i=2}^m c_i
  e_i\right). \]
This gives us a path in $\GL_m$ from the identity to
$\left[\frac{1}{c_1} a_1,e_2,\dots,e_m\right]$:
\[ [0,1]\ni t \mapsto \left[ e_1 + t\left(\sum_{i=2}^m
    \frac{c_i}{c_1}e_i\right),e_2,\dots,e_m \right] \in \GL_m. \]
This process can be continued, since we can rewrite the similar
expression for $a_2$ as
\[ \frac{1}{d_2}a_2 = e_2 + \frac{1}{d_2}\left(\sum_{i=3}^m d_i
  e_i\right) + \frac{d_1}{d_2c_1}a_1, \]
and get a path
\[ t\mapsto \left[\frac{1}{c_1}a_1,e_2+t\left(\sum_{i=3}^m
    \frac{d_i}{d_2}e_i +
    \frac{d_1}{d_2c_1}a_1\right),e_3,\dots,e_m\right] \]
that starts where the other path stopped. If we continue through the
vectors in $A$, we will eventually end at a point,
\[ [\l_i e_i,\dots,\l_me_m,\l_1e_1,\dots,\l_{i-1}e_{i-1}]\in \GL_m, \]
which is diagonal, except the columns may have been
permuted if $n$ is not divisible by $m$. Choosing a fixed path in
$\GL_m$ from $[e_i,\dots,e_{i-1}]$ to $\Id$, scaling the columns and
appending it to all the other paths then gives a path $ \g_A$ with
\[  \g_A(0) = \Id, \qquad \g_A(1) =
[\l_1e_1,\dots,\l_me_m]. \]

If we consider this as a path in $\GLB{m}$ we
get a loop, since $\g_A(1)\in\B_m$ is a diagonal matrix. The
assignment
\[ f(A) = \g_A \in \Omega(\GLB{m}) \]
defines a continuous function from $Y_{m,n}$ to the loop space. Note
that \fixme{Argumenter for kontinuitet?}
this function descends to the quotient space $X_{m,n}=Y_{m,n}/\T^n$, since
elements of $Y_{m,n}$ that are equivalent under the $\mathrm{T}^n$-action gives
paths in $\GL_m$ that are equivalent under the $\mathrm{B}_m$-action.

The goal of this chapter is to prove the following theorem:

\begin{theorem}
  The map $f : X_{m,n} \to \Omega$ induces a map
  \[ f : X_{m,\infty} \to \Omega. \]
  This is a homotopy equivalence.
\end{theorem}

Since the space $\Omega$ is relatively well-studied, this gives
information about $X_{m,\infty}$. In particular, it allows us to
calculate the homology and cohomology of this space.

We want to show that the stabilization map
\begin{align*}
  s : X_{m,n} &\to X_{m,n+1} \\
  A\cdot \T^n = (a_1,\dots,a_n)\cdot \T^n 
              &\mapsto \left(La_1,\dots,La_n,
    \begin{pmatrix}
      1 \\
      \vdots\\
      1
    \end{pmatrix} 
\right) \cdot \T^n = (LA,Le_1) \cdot \T^n
\end{align*}
does not affect the map $f$, or more precisely that $f$ induces a
well-defined map on $X_{m,\infty}$. To do this we will work with the
subsequence $X_{m,nm}$ of the sequence $X_{m,n}$, as the two limits
are isomorphic. \fixme{Argumenter for at grænserne er isomorfe?} To
ease the notation, we will leave out the coset specifier $\cdot \T^n$
when discussing elements of $X_{m,n}$ for the remainder of this
chapter.

We will show that the diagram
\[ \xymatrix{ X_{m,nm} \ar[rr]^{s^m}\ar[dr]_f & & X_{m,(n+1)m}
  \ar[dl]^f \\
  & \Omega & } \]
commutes up to homotopy. This is done by showing that the map $s^m$ is
homotopic to the map $\widetilde s$ that inserts the identity on the
last $m$ columns of $A \in X_{m,nm}$:
\[ \widetilde s (a_1,\dots,a_{nm}) =
(a_1,\dots,a_{nm},e_1,\dots,e_m). \]
If we connect $\Id$ and $L^{-m}$ by a path $\g$ in the space of
invertible lower triangular matrices, we get a homotopy
\[ s_t(A) = (\g(t)L^mA,
\g(t)L^me_1,\g(t)L^{m-1}e_1,\dots,\g(t)Le_1) \]
which goes from $s_0 = s^m$ to the map
\[ s_1 : A \mapsto (A,e_1,L^{-1}e_1,\dots,L^{-m+1}e_1). \]
The matrix $T = [e_1,L^{-1}e_1,\dots,L^{-m+1}e_1]$ is upper
triangular with entries given by the formula:
\[ T_{ij} = (L^{-j+1}e_1)_i = (-1)^{i-1}\binom{j-1}{i-1} \quad
\text{when } i \leq j. \]
This is proven by considering negative powers of $L$ and showing the
identity:
\[ (L^{-j+1})_{ik} = (-1)^{i-k}\binom{j-1}{i-k} \quad \text{when } i
\geq k. \] \fixme{Gør dette}
This is automatically true when $j$ is one, so assume it is true for
a given $j$. We will then show it for $j+1$. We are considering the
entry
\[ (L^{-(j+1)+1})_{ik} = (L^{-j})_{ik} = (L^{-1}\cdot L^{-j+1})_{ik} =
\sum_{l=1}^m (L^{-1})_{il} \cdot (L^{-j+1})_{lk}. \]
But if we use that $L^{-1}$ has entries
\[ (L^{-1})_{il} =
\begin{cases}
  -1 & l = i-1 \\
  1 & l = i  \\
  0 & \text{otherwise}
\end{cases}, \] \fixme{Skriv dette når $L^{-1}$ introduceres}
this simplifies to
\begin{align*}
  (L^{-(j+1)+1})_{ik} &= (-1)\cdot(-1)^{i-1-k}\binom{j-1}{i-1-k}
                        + (-1)^{i-k}\binom{j-1}{i-k} \\
                      &= (-1)^{i-k}\cdot \left( \binom{j-1}{i-k-1} +
                        \binom{j-1}{i-k} \right)\\
                      &= (-1)^{i-k} \binom{j}{i-k} \\
                      &= (-1)^{i-k} \binom{(j+1)-1}{i-k}
\end{align*}
as desired. Since we are considering only the first column, we get
\[ T_{ij} = (L^{-j+1}e_1)_i = (L^{-j+1})_{i1} =
(-1)^{i-1}\binom{j-1}{i-1}. \]
As $T$ is an upper triangular matrix, it can be homotoped to
the identity without affecting the linear independence of the
columns, and the maps $s^m$ and $\widetilde s$ are therefore
homotopic. The composition
$f\circ \widetilde s$ is the same as $f$, except that the loop
$f\circ\widetilde s$ is stationary at the identity for a while before
it terminates. Putting it all together shows the following:
\begin{lemma}
  The diagram
  \[ \xymatrix{ X_{m,nm} \ar[rr]^{s^m}\ar[dr]_{f} & & \ar[dl]^f
    X_{m,(n+1)m} \\
    & \Omega & } \]
  commutes up to homotopy and $f$ induces a map:
  \[ f : X_{m,\infty} \to \Omega. \]
\end{lemma}

\section{Homotopy equivalence}

To show that the map $f : X_{m,\infty} \to \Omega$
gives an isomorphism of homotopy groups, we will again follow
\cite{milnor}. For each element $U$ of $\SUT{m}$, we can find a 
geodesically convex open ball centered on $U$. Since $\SUT{m}$
is compact, there is some $\e > 0$ such that any ball of radius $\e$
or less is geodesically convex. By the same argument, we can also
choose $\e$ so small that for any two points $A= [a_1,\dots, a_m]$ and
$B=[b_1,\dots,b_m]$ in $\SUT{m}$, with the distance from $A$ to $B$ less
than $\e$, any $m$ subsequent vectors in the sequence
$(a_1,\dots,a_m,b_1,\dots,b_m)$ are linearly
independent. \fixme{Argumenter for dette?}
With this, we get an increasing sequence of open sets in $\Omega$,
defined by requiring the loops to be piecewise contained in a ball of
radius $\e$:
\[ \Omega_k = \set{\g \in \Omega \delim \Big\vert\delim
  \g_{\left|\left[ \frac{j-1}{2^k},\frac{j}{2^k} \right]\right.} \text{ is
  contained in an $\e$-ball $\mathrm{B}_\e$} }. \]
By the choice of $\e$ made above, we can define a map
\begin{align*}
  \varphi_k : \Omega_k &\to X_{m,(2^k-1)m} \\
  \g &\mapsto \left[ \g\left(\frac{1}{2^k}\right),
    \g\left(\frac{2}{2^k}\right), \dots,
    \g\left(\frac{2^k-1}{2^k}\right) \right].
\end{align*}
Note that $\g(t)$ is an element of $\SUT{m}$, so the notation
$[\g(t),\dots]$ is shorthand for $[\g(t)e_1,\dots,\g(t)e_m,\dots]$.
The map is continuous since two paths being close in $\Omega_k$ means
that the maximum distance between them is small. But this translates
to the vectors of the image being close together. The maps $\varphi_k$
and $\varphi_{k+1}$ are related in the following fashion:

\begin{lemma}
  \label{lem:com}
  The following diagram commutes up to homotopy:
  \[ \xymatrix{\Omega_k \ar@{^(->}[r] \ar[d]^{\varphi_k} & \Omega_{k+1}
  \ar[d]^{\varphi_{k+1}} \\
  X_{m,(2^k-1)m} \ar[r]^{s} & X_{m,(2^{k+1}-1)m}. } \]
\end{lemma}
\begin{proof}
  The proof is by direct computation.
  \begin{align*}
    \varphi_{k+1} &= \left( \g \mapsto \left[
        \g\left(\frac{1}{2^{k+1}}\right),
        \g\left(\frac{2}{2^{k+1}}\right), \dots,
        \g\left(\frac{2^{k+1}-1}{2^{k+1}}\right) \right]\right) \\
    &\simeq \left(\g \mapsto \left[ \g\left(\frac{2}{2^{k+1}}\right),
      \g\left(\frac{2}{2^{k+1}}\right),
      \g\left(\frac{4}{2^{k+1}}\right), \dots,
      \g\left(\frac{2^{k+1}}{2^{k+1}}\right) \right]\right) \\
    &\simeq \left(\g\mapsto \left[ \g\left(\frac{1}{2^{k}}\right),
      \g\left(\frac{2}{2^{k}}\right), \dots,
      \g\left(\frac{2^k-1}{2^k}\right), \Id,\dots,\Id \right] \right)\\
    &= s\circ \varphi_k.
  \end{align*}
  The first homotopy is given by sliding along the minimal geodesic
  between the points and gives an allowed path by our choice of
  $\e$. The second homotopy is somewhat similar. If $(a_1,\dots,a_m)$
  and $(b_1,\dots,b_m)$ are two elements of $\SUT{m}$ and we have an
  element of $X_{m,n}$ of the form
  \[ \left[ \dots, a_1,\dots,a_m,a_1,\dots,a_m,b_1,\dots,b_m,\dots
  \right], \]
  then we can change the middle columns one by one, as we did in
  the construction of $f$, and get a path in $X_{m,n}$ that changes
  the middle from $[a_1,\dots,a_m]$ to $[b_1,\dots,b_m]$:
  \begin{align*}
    \left[a_1,\dots,a_m,a_1,a_2,\dots,a_m,b_1,\dots,b_m \right] 
    &\leadsto
    \left[a_1,\dots,a_m,b_1,a_2,\dots,a_m,b_1,\dots,b_m\right] \\ 
    &\leadsto
    \left[a_1,\dots,a_m,b_1,b_2,\dots,a_m,b_1,\dots,b_m\right] \\  
    &\leadsto
    \left[a_1,\dots,a_m,b_1,b_2,\dots,b_m,b_1,\dots,b_m\right].  
  \end{align*}
  This allows us to take all the repetitions
  $\left[\dots,\g(t),\g(t),\dots\right]$ and move them to the end as
  diagonal matrices, which results in the stabilization map $s$.
\end{proof}

The above lemma shows that the maps $\varphi_k$ fit together to define
a map from $\Omega$ to $X_{m,\infty}$.
The idea is to use the maps $\varphi_k$ as homotopy inverses of the
loop space map $f$. We are now ready to prove the result.

\begin{theorem}
  The map $f : X_{m,\infty} \to \Omega$ is a weak homotopy
  equivalence, i.e. the map
  \[f_* : \pi_i(X_{m,\infty}) \to \pi_i(\Omega)\]
  is an isomorphism for all $i$.
\end{theorem}

\fixme{Enten en bemærkning om at $\Omega$ og $X$ begge er CW og dermed
  svag h.æ. => h.æ., eller skriv svag i ``hovedsætningen''. Er $X$ CW?
  Kan vi vise det let?}

\begin{proof}
  We will start with surjectivity. Let
  \[ g : W \to \Omega \]
  be a finite cell complex in $\Omega$. We want to show that we can
  construct a map $\widetilde g$ from $W$ to $X_{m,\infty}$ such that
  the composition $f\circ \widetilde g$ is homotopic to $g$.
  
  The image of $g$ is compact and hence contained in $\Omega_k$ for
  some $k$, since $\set{\Omega_k}_{k\in\N}$ is an open cover of
  $\Omega$. Consider the map 
  \[ \widetilde g = \varphi_k\circ g : W \to X_{m,(2^k-1)m}. \]
  Evaluating at a point $w\in W$ gives
  \[ \widetilde g(w) = \varphi_k(\g_w) = \left[
    \g_w\left(\frac{1}{2^k}\right),\dots,
    \g_w\left(\frac{2^k-1}{2^k}\right) \right]. \]
  Evaluating $f$ on this gives a path connecting the points
  $\g_w\left(\frac{j}{2^k}\right)$. But by our choice of $\e$, any
  such path is homotopic to the path connecting these points with
  minimal geodesics, so $f\circ\widetilde g \simeq g$.

  To show injectivity, let
  \[ g : W \to X_{m,\infty}\]
  be a finite cell complex with $f\circ g$ null-homotopic. We want to
  show that $g$ is null-homotopic in $X_{m,\infty}$. In the
  direct limit, a compact set is contained in the image of $X_{m,n}$
  for some $n$ (see \cite[Proposition~A.1]{hatcher} or
  \cite[Chapter~9.4]{may}), so $g$ is represented by a map:
  \[ g : W \to X_{m,n}. \]
  We will work with this representative. The map $f\circ
  g$ is null-homotopic in $\Omega$ and the image is compact, so it is
  contained in $\Omega_k$ for some $k$. By applying $\varphi_k$, we
  get a function
  \[ \varphi_k\circ f\circ g : W \to X_{m,(2^k-1)m}. \]
  Since $f\circ g$ is null-homotopic, this function is
  null-homotopic. But by the same argument that was used in
  Lemma \ref{lem:com}, the map
  \[ \varphi_k(f(g(w))) = \left[
    \g_w\left(\frac{1}{2^k}\right),\dots,\g_w\left(\frac{2^k-1}{2^k}
    \right)\right] \] 
  is homotopic to the stabilization map
  \[ s(g(w)) = \left[g(w),\Id,\dots,\Id\right], \]
  with the homotopy given by sliding along the path $\g_w=f(g(w))$.
  So the map $g$ can be stabilized to a null-homotopic map, so it must
  be null-homotopic in $X_{m,\infty}$. This shows that $f_*$ is
  injective on homotopy groups.
\end{proof}

Since $f$ is a weak homotopy equivalence, it follows that it induces
an isomorphism on homology and cohomology groups. The Pontrjagin
homology ring of the loop space, with product given by loop
concatenation, has been calculated in \cite[Theorem~4.1]{grbic} as
\begin{align*}
  H_*(\Omega(\SUT{m})) \cong T(x_1,\dots,x_{m-1}) \otimes
  \Z[y_1,\dots,y_{m-1}],
\end{align*}
where $T(x_1,\dots,x_{m-1})$ is the tensor algebra generated by
$x_1,\dots,x_{m-1}$. The product has the relation $x_k^2 =
x_px_q+x_qx_p = 2y_1$ for $1\leq k,p,q\leq m-1$ and $p\neq q$. The
degree of $x_i$ is 1 and the degree of $y_i$
is $2i$. By the above result, this is also the homology of
$X_{m,\infty}$. The groups are finitely generated and torsion free, so
by the Universal Coefficient Theorem the homology and cohomology
groups are isomorphic as abelian groups.



\section{Other coefficients}
\label{sec:coef}
\fixme{Andre koefficienter, kan nok gøres ret kort her. Skal ca. bare
  bemærke at alting virker}


\section{Cohomological stability}
\label{sec:costa}
Since we now know the space $X_{m,\infty}(\sigma)$ its cohomology, it
would be nice to relate this to the cohomology of
$X_{m,n}(\sigma)$. To do
this, we specialise to the case $m = 3$ in everything that
follows. First, we need an offset that depends on the permutation
we are considering:
\begin{definition}
  Define a function $\lambda$ on the symmetric group of three elements
  by
  \begin{align*}
    \lambda(\sigma) &=
                      \begin{cases}
                        0 & \sigma = \Id \\
                        1 & \sigma \in \set{(1\,2),(2\,3),(1\,2\,3)}
                        \\
                        2 & \sigma \in \set{(1\,3\,2), (1\,3)}
                      \end{cases}
  \end{align*}
\end{definition}

This section will prove the following theorem:
\begin{theorem}
  If $n \geq k + \lambda(\sigma)$, the stabilisation map
  \[ L_{n}^* : H^k(X_{3,\infty}(\sigma)) \to H^k(X_{3,n}(\sigma)) \]
  is an isomorphism. \fixme{Snak om stabilisering forfra vs. bagfra, L
  vs. s}
\end{theorem}

The proof will proceed by induction on $k$. Our induction hypothesis
is divided into two parts, were we require the stabilisation to be an
isomorphism when $n$ is greater than $k+\lambda(\sigma)$, and
injective when $n$ is $k+\lambda(\sigma)-1$. The exact statement we
want to prove is the following:
\begin{hypothesis}[$I_{k-1}$]
  For all $r \leq k-1$ and all $\sigma \in S_3$, the following two
  statements hold:
  \begin{itemize}
  \item The map $L_{n}^* : H^r(X_{3,\infty}(\sigma)) \to
    H^r(X_{3,n}(\sigma))$ is an isomorphism for $n \geq r +
    \lambda(\sigma)$.
  \item The map $L_{n}^* : H^r(X_{3,\infty}(\sigma)) \to
    H^r(X_{3,n}(\sigma))$ is injective for $n \geq r +\lambda(\sigma)
    -1$.
  \end{itemize}
\end{hypothesis}

By direct calculation, the
groups $H^0(X_{3,n}(\sigma))$ are all $\Z$ when $n \geq
\lambda(\sigma)$, so the isomorphism in case $k = 0$ is automatically
true. \fixme{måske bemærke noget med at injektivitet ikke giver mening
  her}

We now assume the hypothesis $I_{k-1}$ and want to prove $I_k$. Before
we start, we consider the spectral sequence for $X_{3,n+1}(\sigma)$
and observe how the different permutations that show up relate to each
other. This is summarised in the following table:

\[
\begin{array}{|c|cccc|}
  \hline
  \sigma & \multicolumn{1}{c|}{\sigma_\emptyset} 
  & \multicolumn{1}{c|}{\sigma_{\set{1}}} 
  & \multicolumn{1}{c|}{\sigma_{\set{2}}} & \sigma_{\set{1,2}} \\
  \hline
  \Id & \Id & (1\,2) & (2\,3) & (1\,3\,2) \\
  (1\,2) & \Id & \Id & (1\,2\,3) & (1\,3) \\
  (2\,3) & \Id & \Id & (1\,2) & (1\,2) \\
  (1\,2\,3) & \Id & \Id & \Id & \Id \\
  (1\,3\,2) & (2\,3) & (1\,2\,3) & (2\,3) & (1\,2\,3) \\
  (1\,3) & (2\, 3) & (2\,3) & (2\,3) & (2\,3) \\
  \hline
\end{array}
\] \fixme{Leg med stregerne og tjek at det faktisk er rigtigt hvad der
  står...}
By considering the table, we can see the following inequalities for
the function $\lambda$:
\begin{align*}
  \lambda(\sigma) &\geq \lambda(\sigma_{\emptyset}), \\
  \lambda(\sigma) &\geq \lambda(\sigma_{\set{1}}) - 1, \\
  \lambda(\sigma) &\geq \lambda(\sigma_{\set{2}}) - 1, \\
  \lambda(\sigma) &\geq \lambda(\sigma_{\set{1,2}}) - 2. \\
\end{align*}
These will be used during the proof.

\subsection{Injectivity}

We will start by proving injectivity. Consider a permutation $\sigma$
and $n \geq k + \lambda(\sigma) - 1$. We will start by proving that
the stabilisation map
\[ s^* : H^k(X_{3,n+1}(\sigma)) \to
H^k(X_{3,n}(\sigma_{\emptyset})) \]
is injective. Consider the spectral sequence for computing
$H^k(X_{3,n+1}(\sigma))$. Since we choose to stabilise from the left,
as discussed in \fixme{Diskuter dette...Hvor?}, the map $L_n$
becomes a map of spectral sequences and respects the stabilisation map
$s$. The
following is an image of the first page of the spectral sequence, with
... denoting that $L_n^*$ is an isomorphism and ... denoting
that it is injective. \fixme{Tegn og indsæt her}

If the map 
\[ s^* : H^k(X_{3,n+1}(\sigma)) \to
H^k(X_{3,n}(\sigma_{\emptyset})) \]
is not injective, there is an element $\alpha$ of the kernel of the
boundary map that is also in the group 
\[ H^{k-1}(X_{3,n}(\sigma_{\emptyset})) \otimes t_1 \oplus 
H^{k-1}(X_{3,n}(\sigma_{\emptyset})) \otimes t_2 \oplus 
H^{k-2}(X_{3,n}(\sigma_{\emptyset})) \otimes t_1 t_2. \]
But on these last groups, the stabilisation map is an isomorphism so
we get an element $\widehat{\alpha}$ in the corresponding group in the
spectral sequence for $X_{3,\infty}(\sigma)$. Since the stabilisation
maps are maps of spectral sequences, we have
\[ 0 = d_1(L_{n}(\widehat{\alpha})) = L_n(d_1(\widehat{\alpha})) \]
But the right hand side is the composition of $L_n$, which we know by
the induction hypothesis to be injective, and the boundary map in the
spectral sequence for the infinite spaces, which we know to be
injective on the groups we are considering. So $\widehat{\alpha}$ must
be zero, and hence so is $\alpha$. This shows that $s^*$ is injective
when $n \geq k + \lambda(\sigma) - 1$. Now we consider the different
choices of $\sigma$.

\subsubsection{$\sigma = \Id$}

We know that the cohomology of $X_{3,\infty}$ is the inverse
limit of the cohomology groups of $X_{3,n}$ using the stabilisation
map. \fixme{Sørg for at vi ved det, eller skriv det her. Henvis til
  May}
When considered in this way, the map $L_n$ is just the projection on
the corresponding factor. But if we project and get zero, we know that
the only possible lift to the next group is zero, since $s^*$ is
injective. Continuing like this, we see that the only element of
$H^k(X_{3,\infty})$ that can project to zero is the zero element, and
hence $L_n$ is injective in this case.

\subsubsection{$\lambda(\sigma) = 1$}

If $\sigma \in \set{(1\,2),(2\,3),(1\,2\,3)}$, we know that
$\sigma_{\emptyset} = \Id$ and that $n \geq k + \lambda(\sigma) - 1 =
k$, giving $n-1\geq k-1$. Consider the following commutative diagram:
\[ \xymatrix{ H^k(X_{3,\infty}(\sigma)) \ar[r]^{s^*}_{\cong} \ar[d]^{L_n} &
  H^k(X_{3,\infty}) \ar[d]^{L_{n-1}} \\
  H^k(X_{3,n}(\sigma)) \ar[r]^{s^*} & H^k(X_{3,n-1}
} \]
By the previous case, the map on the right is injective and so the
composition $s^* \circ L_{n-1}$ is as well. Since the map commutes,
this means that the left vertical map, $L_n$ must be injective as
well, as desired.

\subsubsection{$\lambda(\sigma) = 2$}

The last case has $\sigma_{\emptyset} = (2\,3)$ and proceeds exactly
as the previous case, using the inectivity that has already been
proved for this permutation.

\subsection{Surjectivity}




% \section{Future work}

% In the above, we showed that the space 
% \[ X_{m,\infty} = \varinjlim_n X_{m,n}/\mathrm{T}^n \]
% is equivalent to the loop space $\Omega(\SUT{m})$ and used this to
% calculate the homology and cohomology groups of the limit. The
% current aim of the project is to work with the spaces $X_{m,n}$
% directly and see what we can prove about these, in particular by using
% the stabilization map
% \[ s : X_{m,n} \to X_{m,n+1} \]
% The conjecture at the moment is the following:
% \begin{conjecture}
%   The induced maps on cohomology
%   \[ s^* : H^*(X_{m,n+1}) \to H^*(X_{m,n}) \]
%   are always surjective and the induced maps on homology
%   \[ s_* : H_*(X_{m,n}) \to H_*(X_{m,n+1}) \]
%   are always injective.
% \end{conjecture}
% The idea for a proof would be to use a spectral sequence argument like
% the one applied to $X_{3,3}$ by defining a filtration that is easier
% to work with and then proceeding by induction, using that the maps
% $s^*$ and $s_*$ induce maps of the spectral sequences.

% Another thing to consider is modifying the space $X_{m,n}$. One possible
% change could be to replace $X_{m,n}$ by a union:
% \[ X'_{m,n} = \bigcup_{A\in\GL_m} X_{m,n}(A,A) \]
% This would allow all sequences of vectors that start and end at the
% same matrix, and one would expect the space $X'_{m,\infty}$ to be
% equivalent to the free loop space, consisting of all
% maps from the circle $S^1$ to $\SUT{m}$.

% There is also the option of replacing the complex numbers $\C$ with
% some other space, like the reals or quaternions. The fact that $\R^*$
% is not connected would most likely make the space $X_{m,n}^{\R}$ quite
% different, but possibly the space $X_{m,n}^{\mathbb{H}}$ would behave in a
% similar manner to $X_{m,n}$, with some shifts in dimension.


%%% Local Variables: 
%%% mode: latex
%%% TeX-master: "main"
%%% End: 
 % Introducer også \mathbb{R} og \mathbb{H}
                  % tilfældene her, og vis sætningen for disse. Dette
                  % skal måske rykkes op før de lidt kedelige
                  % kohomologi-udregninger. Det er nok også her man
                  % skal snakke om stabilisering helt generelt. Og man
                  % kan henvise til H^1-halløjet for at vise hvorfor
                  % det er smart eller lignende.

%\include{tre} % Specialiser til X_{3,n}-tilfældet og snak om
               % stabilisering etc. Jeg er lidt usikker på præcis hvad
               % der skal indgå her.


\chapter{Future work}

I will end with a brief discussion of possible future work that could
be done to study the spaces $X_{m,n}(\sigma)$ and their relation to
the loop space.

\section{Stabilisation}

An obvious first place to start would be to prove Theorem
\ref{thm:kostab} for all $m$, giving something like the following.

\begin{conjecture}
  For each $m$, there is a function $\lambda$ from the symmetric
  group $S_m$ to the integers, such that the stabilisation map
  \[ L_n^* : H^k(X_{m,\infty}(\sigma)) \to H^k(X_{m,n}(\sigma)) \]
  is an isomorphism for $n \geq k + \lambda(\sigma)$.
\end{conjecture}

The exact definition of $\lambda$ is probably one of the hard parts of
this conjecture. It will need to satisfy something similar to the
inequalities stated in the previous chapter, which would probably look
like
\begin{align*}
  \lambda(\sigma) &\geq \lambda(\sigma_{\emptyset}), \\
  \lambda(\sigma) &\geq \lambda(\sigma_{\set{1}}) - 1,\\
  &\vdots \\
  \lambda(\sigma) &\geq \lambda(\sigma_{\set{m}}) - 1,\\
  &\vdots \\
  \lambda(\sigma) &\geq \lambda(\sigma_{\set{1,\dots,m}}) - m + 1.
\end{align*} \fixme{Overvej dette igen. Og burde nok kun være
  1,...,m-1 i indekset, men kan ikke lige få det formuleret på en
  ordentlig måde}
By considering the function defined on $S_3$ in the previous chapter,
it seems like the definition should depend on $\sigma(1)$, but it may
also be related to the Bruhat ordering on $S_m$, which means it
depends on the minimum number of simple transpositions required to
write $\sigma$ as a product.

The diagram obtained by considering $L_n^*$ as a map of spectral
sequences,
\[ L_n^* : E_*^{*,*}(\infty,\sigma) \to E_*^{*,*}(n+1,\sigma), \]
also played a role in the proof. The conjecture above will probably
also require a study of this diagram in more generality than was
considered here. There is a risk that the spectral sequence for
$X_{m,\infty}(\sigma)$ does not collapse as readily for high $m$ as it
does for $m=3$, in particular because increasing $m$ increases the
number of non-zero columns on the first page of the spectral sequence,
$E_1^{*,*}(\infty,\sigma)$. Likewise, the maps induced from $L_n^*$
will also have to be well-behaved on the different columns of the
spectral sequence if the proof is to carry over without too many
changes. This, in particular, is where the choice of $\lambda$ plays
a crucial role.

\subsection{Inverse stabilisation}

During the study of these spaces, some computer programs were written
to simplify some computations. By continuing to work with these, and
applying some conjectures about the structure of the spectral
sequence, we arrived at a conjecture that relates the groups
$H^k(X_{3,n}(\sigma))$ and $H^k(X_{3,\infty}(\sigma)$ for $k$ larger
than $n$. This is the opposite situation to what was considered in
Theorem \ref{thm:kostab}, so a proof of this would mean that we
would know almost all the cohomology groups of $X_{3,n}(\sigma)$,
without doing any explicit calculations. The conjecture is very loose
at the moment, but would probably look something like the following.
\begin{conjecture}
  Blah \fixme{Find ud af hvad der skal være her}
\end{conjecture}

And like above, if this could be established it would seem natural to
consider $m > 3$ and try to prove a similar theorem for
$X_{m,n}(\sigma)$.

\subsection{Other coefficients}

In a similar fashion, it seems very natural to try and prove the
stabilisation theorem for the other possible
coefficient-groups. Nothing in the proof seems to rely on the
coefficients being complex numbers, so the proof should be almost
automatic if a suitable $\lambda$ could be defined.

\section{Generalised spaces}

Another avenue of study could be to consider all the spaces
$X_{m,n}(\sigma)$ at the same time:
\[ \widehat{X}_{m,n} = \bigcup_{A \in \GL_m(\C)} X_{m,n}(A,A). \]
Since the space $X_{m,n}$ approximates the loop group when $n$ is
large, by having a sequence of vectors map to a loop at the identity,
one could expect $\widehat{X}_{m,n}$ to approximate the free loop
group of all maps from the circle to $\GL_m(\C)$, by mapping a
sequence of vectors starting and ending at $A$ to a loop at $A$. This
would require a study of similarly defined stabilisation maps for
these spaces and seeing what happens when we take a limit. This
is summarised in the following conjecture.

\begin{conjecture}
  There are suitable stabilisations maps
  \[ s : \widehat{X}_{m,n} \to \widehat{X}_{m,n+1}, \]
  so that the space
  \[ \widehat{X}_{m,\infty} = \varinjlim_n \widehat{X}_{m,n} \]
  is homotopy equivalent to the free loop space
  \[ \mathrm{L}(\GL_m(\C)) = \set{ \gamma : S^1 \to \GL_m(\C) }. \]
\end{conjecture}

Like the spaces $X_{m,n}$ considered here, this would give another way
to study a complicated space, $\mathrm{L}(\GL_m)$, by working with
simpler versions that consist of a finite sequence of points that we
connect with nice paths.

\fixme{En eller anden form for afslutning...}

% Slå sammen med sidste kapitel?

% Se bagside af blok. Blah blah blah
% Formodning om at s^* er surjektiv.
% Erstat \C med \R eller \H
% Se om rummene A....A approksimerer det frie løkkerum

%%% Local Variables: 
%%% mode: latex
%%% TeX-master: "main"
%%% End: 
 % Noget om de ting vi ikke nåede og generel
                    % spekulation. Hvad kan skidtet eventuelt bruges
                    % til, fremtidigt arbejde (hah) og lignende.

% Skal der noget som helst med om det andet jeg lavede med pi_1?
% Sikkert ikke, men snak lige med Marcel.

\appendix

\chapter{Sage code} 
\label{appendix:sage}

The Sage code that was used for computations, referenced in Chapter
\ref{chap:udregninger}, is included below. Note that a red arrow at
the beginning of a line of code denotes a linebreak that was inserted
to make the code fit on the page and is not there in the original
code.

\section{The boundary map}
\label{sec:sageboundary}

The code below computes the kernel and image of the first boundary map
in the spectral sequence for $Y_{3,3}$, as described in Section
\ref{sec:rand}. It does so by creating a ring with the correct
generators and relations and encoding the boundary map in two steps,
firstly using the inclusion of the relevant spaces and then using the
boundary map from the pair $(D,D-0)$. It the splits the cohomology
groups appearing in the spectral sequence into ranks and computes
kernel and image for the boundary. To get a readable output, load the
Sage file and call either the function \lstinline{pretty_print} for
the full $E_2$-page of $Y_{3,3}$ or \lstinline{print_mod} for the
cohomology of $X_{3,3}$.

\fixme{Skriv kommentarer i koden}
\lstinputlisting[language=Python]{./sage/rand.sage}


\section{Orientation of intersections}
\label{sec:sageorientation}

The code below computes the orientation of the points of intersection
of the inclusion maps, as described in Section
\ref{sec:orientation}. Note that the code is old, and hence uses a
slightly different naming convention than is otherwise used in this
thesis. In particular, the space $X_{001}$ in the code corresponds to
the space $Y_{\set{1,2}}$ in the thesis, $X_{101}$ corresponds to
$Y_{\set{1}}$, $X_{011}$ to $Y_{\set{2}}$ and $X_{111}$ to $Y_{\emptyset}$.

\lstinputlisting[language=Python]{./sage/orientation.sage}


%%% Local Variables: 
%%% mode: latex
%%% TeX-master: "main"
%%% End: 


\backmatter
\nocite{*}
\printbibliography{}

\end{document}
